% =====================================================================================================
%  پایان‌نامه‌ی کارشناسی: سامانه‌ی بازیابی هوشمند اطلاعات و استخراج خواص فیزیکی مواد با معماری RAG
%  کامپایلر Overleaf: XeLaTeX     |     ساخت محلی: bash .claude/skills/persian-thesis-pdf/scripts/build.sh main.tex
%  تصاویر با \figbox{name.png}{width} درج می‌شوند؛ اگر فایل نباشد کادر جای‌خالی چاپ می‌شود.
%  عددهای مثال‌های حل‌شده و نمودارهای فصل ۵ را scripts/calc_figures.py می‌سازد.
% =====================================================================================================
\documentclass[12pt,a4paper]{report}

\usepackage{amsmath,amssymb}
\usepackage{graphicx}
\usepackage[a4paper,top=2.8cm,bottom=2.6cm,right=2.6cm,left=2.4cm,headheight=16pt]{geometry}
\usepackage{setspace}
\usepackage{array}
\usepackage{booktabs}
\usepackage{multirow}
\usepackage{enumitem}
\usepackage{fancyhdr}
\usepackage{xcolor}
\usepackage{tikz}
\usetikzlibrary{positioning,arrows.meta}
\usepackage[colorlinks=true,linkcolor=black,citecolor=blue,urlcolor=blue]{hyperref}
\usepackage{xepersian}   % باید آخرین بسته باشد

% ---------------------------------------- قلم یاس ----------------------------------------------------
\IfFontExistsTF{BYas.ttf}{%
  \IfFontExistsTF{BYasBd.ttf}%
    {\settextfont[Scale=1.12,BoldFont=BYasBd.ttf]{BYas.ttf}}%
    {\IfFontExistsTF{BYasBold.ttf}%
      {\settextfont[Scale=1.12,BoldFont=BYasBold.ttf]{BYas.ttf}}%
      {\settextfont[Scale=1.12]{BYas.ttf}}}}%
{\IfFontExistsTF{BYas}{\settextfont[Scale=1.12]{BYas}}%
{\IfFontExistsTF{B Yas}{\settextfont[Scale=1.12]{B Yas}}%
{\IfFontExistsTF{XB Yas}{\settextfont[Scale=1.12]{XB Yas}}%
{\IfFontExistsTF{Yas}{\settextfont[Scale=1.12]{Yas}}%
{\IfFontExistsTF{XB Niloofar}{\settextfont[Scale=1.12]{XB Niloofar}}%
{\IfFontExistsTF{Vazirmatn}{\settextfont[Scale=1.05]{Vazirmatn}}%
{\IfFontExistsTF{Nazli}{\settextfont[Scale=1.12]{Nazli}}%
{\settextfont[Scale=1.05]{Amiri}}}}}}}}}

% ---------------------------------------- فاصله‌ی خطوط --------------------------------------------
\setstretch{1.55}
\setlength{\parskip}{3pt}
\setlist{itemsep=4pt,topsep=4pt}
\setlength{\emergencystretch}{3em}

% ---------------------------------------- سرصفحه ---------------------------------------------------
\pagestyle{fancy}
\fancyhf{}
\fancyhead[R]{\small\leftmark}
\fancyfoot[C]{\thepage}
\renewcommand{\headrulewidth}{0.5pt}
\fancypagestyle{plain}{\fancyhf{}\fancyfoot[C]{\thepage}\renewcommand{\headrulewidth}{0pt}}

% ---------------------------------------- فرمان‌های کمکی ------------------------------------------
\newcommand{\sub}[1]{\mathrm{#1}}
\newcommand{\Imag}{\operatorname{Im}}
\newcommand{\Real}{\operatorname{Re}}
\newcommand{\sinc}{\operatorname{sinc}}
\newcommand{\figbox}[2]{%
  \IfFileExists{#1}{\includegraphics[width=#2]{#1}}{%
    \fbox{\parbox[c][\dimexpr#2*42/100\relax][c]{\dimexpr#2-2\fboxsep-2\fboxrule\relax}{\centering\textbf{محل قرارگیری تصویر}\\[4pt]
    فایل تصویر را با نام \lr{\texttt{\detokenize{#1}}} بارگذاری کنید.}}}}
% کادر کد چپ‌به‌راست (برای نمونه‌ی رکورد داده و فرمان‌ها)
\newcommand{\codebox}[1]{\begin{center}\lr{\fbox{\parbox{0.94\textwidth}{\footnotesize\ttfamily\setstretch{1.15}#1}}}\end{center}}

% شکل معماری سامانه (شکل ۳-۱)
\newcommand{\archfigure}{%
\begin{tikzpicture}[
  box/.style={draw,rounded corners=2pt,fill=blue!6,minimum width=3.9cm,minimum height=1.05cm,align=center,font=\small},
  arr/.style={-{Stealth[length=2mm]},thick},
  node distance=0.42cm]
% lane 1
\node[font=\bfseries] (h1) {الف) نمایه‌سازی};
\node[box,below=0.3cm of h1] (a1) {مقاله‌ی \lr{PDF}\\ ۲۸۴۵ سند};
\node[box,below=of a1] (a2) {تبدیل سند\\ \lr{Docling} و \lr{OCR}};
\node[box,below=of a2] (a3) {ترمیم متن\\ حروف یونانی، زیرنویس‌ها};
\node[box,below=of a3] (a4) {بخش‌بندی ساختاری\\ حداکثر ۵۱۲ توکن};
\node[box,below=of a4] (a5) {بردارسازی\\ \lr{bge-base}، ۷۶۸ بعد};
\node[box,below=of a5] (a6) {پایگاه داده‌ی برداری\\ \lr{ChromaDB}، \lr{Qdrant} و \lr{BM25}};
\foreach \i/\j in {a1/a2,a2/a3,a3/a4,a4/a5,a5/a6} \draw[arr] (\i) to (\j);
% lane 2
\node[font=\bfseries,left=0.55cm of h1,xshift=-2.0cm] (h2) {ب) پرسش و پاسخ};
\node[box,below=0.3cm of h2,fill=green!6] (b1) {پرسش فارسی یا انگلیسی\\ ترجمه با واژه‌نامه};
\node[box,below=of b1,fill=green!6] (b2) {بازیابی ترکیبی\\ برداری و \lr{BM25}، ادغام \lr{RRF}};
\node[box,below=of b2,fill=green!6] (b3) {رتبه‌بندی دوم\\ مدل داور روی ۳۰ نامزد};
\node[box,below=of b3,fill=green!6] (b4) {پاسخ با ارجاع\\ مقاله، صفحه، محل روی صفحه};
\node[box,below=of b4,fill=green!6] (b5) {نگهبان «پیدا نشد»\\ نبود شاهد};
\foreach \i/\j in {b1/b2,b2/b3,b3/b4,b4/b5} \draw[arr] (\i) to (\j);
% lane 3
\node[font=\bfseries,left=0.55cm of h2,xshift=-2.0cm] (h3) {ج) داده و ابزار};
\node[box,below=0.3cm of h3,fill=orange!8] (c1) {استخراج ساخت‌یافته\\ \lr{Pydantic} و سه بررسی};
\node[box,below=of c1,fill=orange!8] (c2) {آزمون با عدد مقاله\\ مقدار گزارش‌شده};
\node[box,below=of c2,fill=orange!8] (c3) {پایگاه داده‌ی آزموده\\ $n(\lambda)$، $d_{ij}$ و $dn/dT$};
\node[box,below=of c3,fill=orange!8] (c4) {محاسبه‌گرها\\ تطبیق فاز، \lr{SPDC}، مشاور بلور};
\foreach \i/\j in {c1/c2,c2/c3,c3/c4} \draw[arr] (\i) to (\j);
\end{tikzpicture}}

\begin{document}

% ======================================== صفحه‌ی عنوان ============================================
\begin{titlepage}
\begin{center}
{\large بسمه تعالی}\\[10pt]
\figbox{logo_ui.png}{2.8cm}\\[8pt]
{\Large\textbf{دانشگاه اصفهان}}\\[6pt]
{\large دانشکده‌ی فیزیک}\\[4pt]
{\large گزارش پروژه‌ی کارشناسی فیزیک (۳ واحد)}\\[1.6cm]
\rule{\textwidth}{0.6pt}\\[10pt]
{\LARGE\textbf{طراحی و پیاده‌سازی سامانه‌ی بازیابی هوشمند اطلاعات و استخراج خواص فیزیکی مواد از مقالات علمی}}\\[14pt]
{\large\textbf{با معماری \lr{RAG} و پایگاه داده‌های برداری}}\\[10pt]
\rule{\textwidth}{0.6pt}\\[2.2cm]
\begin{minipage}[t]{0.48\textwidth}
\textbf{پژوهشگر:}\\ امین حسین سدیدی\\[12pt]
\textbf{شماره‌ی دانشجویی:}\\ ۴۰۲۴۲۱۳۰۲۷
\end{minipage}\hfill
\begin{minipage}[t]{0.42\textwidth}
\textbf{استاد راهنما:}\\ دکتر محمدی
\end{minipage}\\[2.4cm]
{\large آبان ۱۴۰۵}
\end{center}
\end{titlepage}

% ======================================== سپاس و اظهارنامه =========================================
\chapter*{سپاسگزاری}
\markboth{سپاسگزاری}{سپاسگزاری}
از استاد راهنمای گرامی، دکتر محمدی، که با راهنمایی‌ها و پیشنهادهای ارزشمندشان در تعریف موضوع و مسیر این پروژه یاری‌ام کردند، صمیمانه سپاسگزارم.

این پروژه بر پایه‌ی ابزارها و داده‌های آزادی بنا شده است که جامعه‌ی علمی در دسترس همگان گذاشته است؛ به‌ویژه پایگاه داده‌ی \lr{refractiveindex.info}، پایگاه \lr{HITRAN}، کتابخانه‌ی \lr{Docling}، مدل برداری \lr{BGE} و پایگاه‌های داده‌ی برداری \lr{ChromaDB} و \lr{Qdrant}.

\chapter*{اظهارنامه}
\markboth{اظهارنامه}{اظهارنامه}
اینجانب امین حسین سدیدی، دانشجوی کارشناسی فیزیک دانشگاه اصفهان به شماره‌ی دانشجویی ۴۰۲۴۲۱۳۰۲۷، گواهی می‌کنم که طراحی و اجرای این پروژه را خود انجام داده‌ام و در استفاده از نتایج پژوهش دیگران، منبع را ذکر کرده‌ام. در نوشتن بخشی از کد و ویرایش متن از دستیارهای هوش مصنوعی استفاده شده است؛ طراحی، تصمیم‌ها، آزمون‌ها و درستی‌آزمایی نتایج بر عهده‌ی اینجانب بوده است.

\vspace{2cm}
\hfill امضا: \rule{4cm}{0.4pt}

% ======================================== چکیده =====================================================
\chapter*{چکیده}
\addcontentsline{toc}{chapter}{چکیده}
\markboth{چکیده}{چکیده}
داده‌های پایه‌ی فیزیک مواد، مانند فرمول پاشندگی ضریب شکست، ضرایب غیرخطی مرتبه‌ی دوم، ضریب ترمواپتیک و پارامترهای خطوط طیفی، در جدول‌ها و معادله‌های هزاران مقاله پراکنده‌اند و مدل‌های زبانی عمومی این مقادیر را اغلب بدون منبع قابل‌اتکا تولید می‌کنند. در این پروژه سامانه‌ای با معماری «تولید تقویت‌شده با بازیابی» (\lr{RAG}) طراحی و پیاده‌سازی شد که مقاله‌های \lr{PDF} را به متن و جدول ساختاریافته تبدیل می‌کند، آن‌ها را در پایگاه داده‌های برداری (\lr{ChromaDB} و \lr{Qdrant}) نمایه می‌کند و به پرسش فارسی یا انگلیسی با ارجاع به مقاله، صفحه و محل روی صفحه پاسخ می‌دهد. بازیابی ترکیبی از جست‌وجوی معنایی با مدل \lr{bge-base-en-v1.5} و جست‌وجوی کلیدواژه‌ای \lr{BM25} است که با روش ادغام معکوس رتبه (\lr{RRF}) یکی می‌شوند. استخراج ساخت‌یافته‌ی خواص با سه بررسی خودکار ضدتوهم همراه است.

سامانه نخست روی طیف‌سنجی مولکولی با مرجع \lr{HITRAN} آزموده و سپس به بلورهای اپتیک خطی و غیرخطی گسترش داده شد. مجموعه‌ی نهایی ۲۸۴۵ سند و ۶۸۱۹۰ بخش قابل‌جست‌وجو دارد. در آزمون \lr{HITRAN}، پاسخ \lr{RAG} با مدل محلی \lr{Qwen3-4B} در ۱۴ پرسش از ۱۵ پرسش مقدار درست را آورد، در حالی که همان مدل بدون منبع در هیچ پرسشی پاسخ درست نداد. در ارزیابی ۲۰۰ پرسش طلایی، \lr{MRR@10} برای جست‌وجوی ترکیبی ۰/۵۱۸ و پس از رتبه‌بندی دوم ۰/۵۹۰ بود. بر این پایه، یک پایگاه داده‌ی فیزیکی آزموده و محاسبه‌گرهای تطبیق فاز و تولید جفت فوتون ساخته شد که خروجی‌شان با مقادیر گزارش‌شده در مقاله‌های اصلی هم‌خوان است؛ برای نمونه، زاویه‌ی تطبیق فاز تولید هماهنگ دوم نوع اول در بلور \lr{BBO} در طول موج ۱/۰۶۴ میکرومتر ۲۲/۸۳ درجه به دست آمد که با مقدار ۲۲/۸ درجه‌ی ادبیات هم‌خوان است. نتیجه‌ی منفی مهم پروژه این است که استخراج کاملاً خودکار ضرایب فرمول با مدل زبانی به دقت لازم نرسید و داده‌ی دقیق فقط با استخراج نیمه‌دستی و آزمون در برابر خود مقاله به دست آمد.

\vspace{6pt}
\noindent\textbf{واژگان کلیدی:} تولید تقویت‌شده با بازیابی، پایگاه داده‌ی برداری، بازیابی ترکیبی، استخراج خواص مواد، ضریب شکست، اپتیک غیرخطی، تطبیق فاز.

\tableofcontents
\listoftables
\listoffigures

% ======================================== فصل ۱ ======================================================
\chapter{مقدمه}

\section{بیان مسئله}
طراحی یک آزمایش فیزیکی معمولاً با چند پرسش عددی ساده شروع می‌شود. ضریب شکست این بلور در طول موج ۱/۰۶۴ میکرومتر چقدر است؟ ضریب غیرخطی مؤثرش چند پیکومتر بر ولت است؟ با چه زاویه‌ای باید بریده شود تا تطبیق فاز رخ دهد؟ ضریب پهن‌شدگی فشاری یک خط جذبی $\mathrm{CO_2}$ چقدر است؟ پاسخ این پرسش‌ها در ادبیات علمی وجود دارد، ولی پیدا کردنش دشوار است، به چهار دلیل:
\begin{itemize}
\item \textbf{پراکندگی:} فرمول سلمایر یک بلور در یک مقاله است، وابستگی دمایی‌اش در مقاله‌ای دیگر و ضرایب غیرخطی‌اش در مقاله‌ی سوم. برای یک بلور پرکاربرد ممکن است ده‌ها اندازه‌گیری مستقل با بازه‌های طول موجی و دمایی متفاوت وجود داشته باشد و انتخاب میان آن‌ها خود نیازمند خواندن دقیق هر مقاله است.
\item \textbf{ناهمگونی:} قرارداد محورها، علامت ضرایب، واحدها و بازه‌ی اعتبار فرمول‌ها از مقاله‌ای به مقاله‌ی دیگر فرق می‌کند. برای نمونه، جابه‌جایی $d_{31}$ و $d_{32}$ در بلورهای \lr{KTP} و \lr{LBO} میان قراردادهای مختلف محور، خطایی رایج است. حتی واحد طول موج در فرمول پاشندگی (میکرومتر یا نانومتر) و تعریف ضریب غیرخطی ($d$ یا $\chi^{(2)}=2d$) همیشه صریح نوشته نمی‌شود.
\item \textbf{قالب سخت‌خوان:} عددها در جدول‌های چندستونه، معادله‌های تصویری و اسکن مقاله‌های دهه‌های ۱۹۶۰ تا ۱۹۹۰ هستند. این قالب‌ها برای خواننده‌ی انسانی قابل‌فهم‌اند، اما ابزارهای خودکار استخراج متن آن‌ها را به‌سادگی خراب می‌کنند.
\item \textbf{هزینه‌ی بالای خطا:} یک عدد اشتباه می‌تواند به خرید بلور نامناسب یا هفته‌ها توقف کار آزمایشگاهی بینجامد. چون بلورهای غیرخطی با زاویه‌ی برش مشخص ساخته می‌شوند، خطای زاویه پس از ساخت قابل‌جبران نیست.
\end{itemize}

مدل‌های زبانی بزرگ (\lr{LLM}) در نگاه نخست راه‌حل این مسئله به نظر می‌رسند، اما در سطح عدد و جدول قابل‌اتکا نیستند. این مدل‌ها پاسخ را از دانش ذخیره‌شده در پارامترهایشان می‌سازند. اگر عدد دقیقی را ندیده باشند یا آن را درست به خاطر نسپرده باشند، عددی باورپذیر ولی نادرست تولید می‌کنند. این پدیده «توهم» (\lr{hallucination}) نام دارد \cite{ji}. این مدل‌ها معمولاً منبع قابل‌بررسی هم نمی‌دهند، پس کاربر راهی برای تشخیص عدد درست از عدد ساختگی ندارد. در داده‌های فیزیکی کم‌تکرار، مانند ضریب پهن‌شدگی یک خط طیفی خاص یا ضریب ترمواپتیک یک بلور کم‌کاربرد، احتمال توهم بیشتر است، زیرا مدل این عدد را در داده‌های آموزشی خود به‌ندرت دیده است.

معماری «تولید تقویت‌شده با بازیابی» (\lr{Retrieval-Augmented Generation} یا \lr{RAG}) برای همین مسئله پیشنهاد شده است \cite{lewis}. در این معماری، پیش از پاسخ، بخش‌های مرتبط از یک مجموعه‌ی سند بازیابی می‌شوند و مدل فقط بر پایه‌ی همان بخش‌ها پاسخ می‌دهد. مسئله‌ی این پروژه، ساختن و سنجیدن چنین سامانه‌ای برای داده‌های فیزیکی است. در این حوزه، معیار موفقیت روانی پاسخ نیست، بلکه درستی عدد و دقت ارجاع آن به منبع است. به بیان دیگر، سامانه‌ای که پاسخ روان ولی بدون ارجاع می‌دهد برای کار فیزیکی ارزشی ندارد، ولی سامانه‌ای که فقط منبع درست را با شماره‌ی صفحه پیدا کند، حتی بدون نوشتن پاسخ، کار پژوهشگر را به‌طور محسوسی کوتاه می‌کند.

\section{اهداف}
هدف کلی پروژه ساختن زنجیره‌ای است که از مقاله‌ی \lr{PDF} شروع شود و به عدد آزموده و قابل‌استفاده در محاسبه‌ی فیزیکی برسد. اهداف جزئی عبارت‌اند از:
\begin{enumerate}
\item طراحی و پیاده‌سازی یک خط لوله‌ی کامل \lr{RAG} برای مقاله‌های فیزیک: تبدیل \lr{PDF}، بخش‌بندی، نمایه در پایگاه داده‌ی برداری، بازیابی و پاسخ.
\item پاسخ به هر پرسش با ارجاع دقیق به مقاله، صفحه و محل روی صفحه، و پاسخ صریح «پیدا نشد» وقتی شاهدی در منابع نیست.
\item استخراج ساخت‌یافته‌ی خواص فیزیکی (مقدار، واحد، شرایط اندازه‌گیری) همراه با بررسی‌های خودکار ضدتوهم.
\item ارزیابی عینی سامانه با پایگاه‌های مرجع مستقل (\lr{HITRAN} و \lr{refractiveindex.info}) به‌جای برچسب‌گذاری دستی حجیم.
\item پشتیبانی از پرسش فارسی برای دانشجویان و پژوهشگران فارسی‌زبان.
\item وصل کردن داده‌ی آزموده به ابزارهای محاسبه‌ی فیزیکی (ضریب شکست، تطبیق فاز، تولید جفت فوتون) تا پاسخ پرسش‌های طراحی آزمایش کمّی و مستند باشد.
\end{enumerate}

\section{پرسش‌های پژوهش}
چهار پرسش زیر مسیر طراحی و ارزیابی را تعیین کردند. هر پرسش به یک تصمیم مهندسی مشخص در سامانه وابسته است و پاسخ هر کدام در فصل ۶ با عددهای فصل ۵ داده می‌شود.
\begin{enumerate}
\item آیا \lr{RAG} در مقایسه با همان مدل زبانی بدون منبع، دقت پاسخ به پرسش‌های عددی فیزیک را بهتر می‌کند؟
\item کدام روش بازیابی (برداری، کلیدواژه‌ای یا ترکیبی) برای متن‌های فیزیکی پر از نماد و فرمول بهتر کار می‌کند؟
\item ترجمه‌ی ماشینی پرسش فارسی به انگلیسی چه اثری بر کیفیت بازیابی دارد؟
\item آیا استخراج کاملاً خودکار ضرایب فرمول‌ها با مدل زبانی به دقت لازم برای کار فیزیکی می‌رسد؟
\end{enumerate}

\section{نوآوری و دامنه‌ی کار}
ابزارهای موجود هر کدام بخشی از این مسئله را حل می‌کنند: \lr{refractiveindex.info} داده‌ی ضریب شکست را با ارجاع دارد \cite{polyanskiy}، نرم‌افزار \lr{SNLO} محاسبه‌ی اپتیک غیرخطی را خوب انجام می‌دهد و دستیارهای پژوهشی مبتنی بر مدل زبانی روی متن مقاله‌ها کار می‌کنند. وجه متمایز این پروژه، زنجیره‌ی کامل از متن مقاله تا محاسبه است. هر عدد به صفحه و جدول مشخصی وصل است، پیش از ورود به پایگاه داده با عدد خود مقاله آزموده می‌شود و سپس در ابزار محاسبه به کار می‌رود. رابط فارسی با واژه‌نامه‌ی ثابت اصطلاحات فیزیک نیز بخشی از کار است.

دامنه‌ی کار دو مرحله داشت. مرحله‌ی نخست، آزمون مفهوم روی پارامترهای طیف‌سنجی گازهای جوّی ($\mathrm{CO_2}$، $\mathrm{CH_4}$ و $\mathrm{H_2O}$) با مرجع \lr{HITRAN} بود \cite{hitran}. این حوزه از آن رو انتخاب شد که پایگاه مرجع آن برای هر پارامتر، مقاله‌ی منبع را ذکر می‌کند و ارزیابی عینی را ممکن می‌سازد. مرحله‌ی دوم، ساخت سامانه‌ی کامل روی حوزه‌ی بلورهای اپتیک خطی و غیرخطی بود. هر دو مجموعه در کد پروژه به‌صورت «مجموعه»‌های (\lr{corpus}) جداگانه نگه داشته می‌شوند تا افزودن یک حوزه‌ی تازه فقط به ساختن یک مجموعه‌ی جدید و فایل پیکربندی آن نیاز داشته باشد.

آنچه خارج از دامنه‌ی این پروژه است نیز باید روشن باشد: سامانه داده‌ی تجربی تازه تولید نمی‌کند، مدل زبانی تازه آموزش نمی‌دهد و جایگزین خواندن مقاله توسط پژوهشگر نیست. هدف، کوتاه کردن مسیر رسیدن به منبع درست و جلوگیری از ورود عدد آزموده‌نشده به محاسبه است.

\section{ساختار گزارش}
فصل ۲ مبانی نظری \lr{RAG}، نمایش برداری، بازیابی کلیدواژه‌ای و مبانی فیزیکی داده‌های هدف را مرور می‌کند. فصل ۳ روش و معماری سامانه را شرح می‌دهد و فصل ۴ جزئیات پیاده‌سازی را. فصل ۵ نتایج ارزیابی را با عددهای اندازه‌گیری‌شده می‌آورد، از جمله نتیجه‌های منفی و دو مثال حل‌شده‌ی فیزیکی. فصل ۶ به بحث، محدودیت‌ها و کارهای آینده می‌پردازد. واژه‌نامه، منشأ عددهای گزارش و دستور اجرای سامانه در پیوست‌ها آمده‌اند.

% ======================================== فصل ۲ ======================================================
\chapter{مبانی نظری و پیشینه}

\section{مدل‌های زبانی و معماری \lr{RAG}}
\subsection{مدل زبانی بزرگ}
مدل زبانی بزرگ یک شبکه‌ی عصبی ترنسفورمر است \cite{vaswani} که با پیش‌بینی واژه‌ی بعدی روی حجم بزرگی از متن آموزش دیده است. متن ورودی ابتدا به دنباله‌ای از «توکن»‌ها (واژه یا پاره‌واژه) شکسته می‌شود و هر توکن به یک بردار نگاشته می‌شود. هسته‌ی ترنسفورمر سازوکار «توجه» است که در آن هر توکن با همه‌ی توکن‌های دیگر دنباله مقایسه می‌شود:
\begin{equation}\label{eq:attention}
\mathrm{Attention}(Q,K,V)=\mathrm{softmax}\!\left(\frac{QK^{T}}{\sqrt{d_k}}\right)V ,
\end{equation}
که در آن $Q$، $K$ و $V$ تصویرهای خطی بردارهای توکن‌ها و $d_k$ بعد این تصویرهاست. مدل در هر گام توزیع احتمال توکن بعدی، $p(y_t\mid y_{<t},x)$، را می‌سازد و پاسخ توکن به توکن تولید می‌شود.

دانش چنین مدلی در پارامترهایش «فشرده» شده است. این دانش پس از آموزش به‌روز نمی‌شود، منبع ندارد و برای عددهای کم‌تکرار ادبیات تخصصی دقیق نیست. دلیل آن روشن است: مدل برای کمینه کردن خطای پیش‌بینی واژه آموزش دیده است، نه برای به خاطر سپردن دقیق رقم‌ها. عددی که در کل داده‌ی آموزشی فقط چند بار آمده، احتمالاً به‌صورت تقریبی یا با عددی شبیه به آن جایگزین می‌شود. افزون بر این، طول ورودی مدل (پنجره‌ی متن) محدود است، پس نمی‌توان هزاران مقاله را یکجا به مدل داد.

\subsection{معماری \lr{RAG}}
در معماری \lr{RAG} \cite{lewis,gao} پاسخ در دو گام ساخته می‌شود. در گام بازیابی، برای پرسش $q$ مجموعه‌ای از $k$ بخش مرتبط $\{c_1,\dots,c_k\}$ از پایگاه اسناد پیدا می‌شود. در گام تولید، مدل زبانی پرسش و این بخش‌ها را با هم می‌بیند و پاسخ را فقط بر پایه‌ی آن‌ها می‌نویسد. در صورت‌بندی احتمالاتی \lr{Lewis} و همکاران، بخش بازیابی‌شده یک متغیر پنهان است:
\begin{equation}\label{eq:rag}
p(y\mid q)\approx\sum_{c\in\mathrm{top}\text{-}k(q)} p_\eta(c\mid q)\,p_\theta(y\mid q,c),
\end{equation}
که در آن $p_\eta$ بازیاب و $p_\theta$ مدل تولیدکننده است. در پیاده‌سازی این پروژه، مانند بیشتر سامانه‌های عملی، به‌جای جمع وزن‌دار، همه‌ی $k$ بخش برتر یکجا در متن ورودی مدل قرار می‌گیرند.

به این ترتیب سه ویژگی به دست می‌آید که برای کار علمی ضروری‌اند: پاسخ قابل‌ردیابی به منبع است، دانش با افزودن سند تازه به‌روز می‌شود و مدل می‌تواند در نبود شاهد بگوید «پیدا نشد». کیفیت پاسخ در این معماری بیش از هر چیز به کیفیت بازیابی وابسته است: اگر بخش درست در میان $k$ بخش بازیابی‌شده نباشد، بهترین مدل زبانی هم پاسخ درست نمی‌دهد. به همین دلیل، بخش بزرگی از ارزیابی این پروژه بر کیفیت بازیابی متمرکز است.

\section{نمایش برداری و پایگاه داده‌ی برداری}
\subsection{رمزگذار برداری}
در جست‌وجوی معنایی، هر بخش متن با یک مدل رمزگذار (\lr{embedding}) به برداری در فضای $d$ بعدی نگاشته می‌شود، به‌گونه‌ای که متن‌های هم‌معنا بردارهای نزدیک داشته باشند \cite{sbert,dpr}. در این پروژه از مدل \lr{bge-base-en-v1.5} استفاده شد \cite{bge} که بردار ۷۶۸ بعدی تولید می‌کند. شباهت پرسش و بخش با شباهت کسینوسی سنجیده می‌شود:
\begin{equation}\label{eq:cosine}
\mathrm{sim}(\mathbf{q},\mathbf{c})=\frac{\mathbf{q}\cdot\mathbf{c}}{\lVert\mathbf{q}\rVert\,\lVert\mathbf{c}\rVert}.
\end{equation}
اگر بردارها از پیش بهنجار شوند ($\lVert\mathbf{q}\rVert=\lVert\mathbf{c}\rVert=1$)، شباهت کسینوسی همان ضرب داخلی است و فاصله‌ی اقلیدسی هم با آن رابطه‌ی یکنوا دارد: $\lVert\mathbf{q}-\mathbf{c}\rVert^2=2-2\,\mathrm{sim}(\mathbf{q},\mathbf{c})$.

رمزگذارهایی مانند \lr{BGE} با یادگیری تقابلی آموزش می‌بینند: برای هر پرسش، یک بخش مرتبط (نمونه‌ی مثبت) و تعدادی بخش نامرتبط (نمونه‌های منفی) در نظر گرفته می‌شود و تابع هزینه شباهت با نمونه‌ی مثبت را نسبت به منفی‌ها بیشینه می‌کند \cite{dpr,bge}:
\begin{equation}\label{eq:infonce}
\mathcal{L}=-\log\frac{\exp\!\left(\mathrm{sim}(\mathbf{q},\mathbf{c}^{+})/\tau\right)}{\exp\!\left(\mathrm{sim}(\mathbf{q},\mathbf{c}^{+})/\tau\right)+\sum_{j}\exp\!\left(\mathrm{sim}(\mathbf{q},\mathbf{c}^{-}_j)/\tau\right)},
\end{equation}
که $\tau$ پارامتر دما است. چون این رمزگذار پرسش و بخش را جداگانه به بردار تبدیل می‌کند (معماری «دورمزگذاری» یا \lr{bi-encoder})، بردار همه‌ی بخش‌ها یک بار و از پیش محاسبه می‌شود و در زمان پرسش فقط بردار خود پرسش لازم است. همین ویژگی جست‌وجو در ده‌ها هزار بخش را در کسری از ثانیه ممکن می‌کند.

\subsection{پایگاه داده‌ی برداری}
پایگاه داده‌ی برداری نرم‌افزاری است که این بردارها را همراه با متن و فراداده‌ی هر بخش ذخیره می‌کند و نزدیک‌ترین بردارها به بردار پرسش را سریع برمی‌گرداند. جست‌وجوی دقیق نیازمند مقایسه‌ی بردار پرسش با همه‌ی $N$ بردار است و هزینه‌ی آن با $N d$ رشد می‌کند. برای میلیون‌ها بردار این روش کند است و از روش‌های تقریبی مانند گراف \lr{HNSW} استفاده می‌شود \cite{hnsw}. در \lr{HNSW} بردارها رأس‌های یک گراف چندلایه‌اند: لایه‌های بالایی پیوندهای بلندبرد و تُنُک دارند و لایه‌ی پایین همه‌ی رأس‌ها را با پیوندهای کوتاه‌برد. جست‌وجو از لایه‌ی بالا شروع می‌شود و با حرکت حریصانه به سمت همسایه‌ی نزدیک‌تر، لایه به لایه پایین می‌آید. هزینه‌ی هر جست‌وجو تقریباً با $\log N$ رشد می‌کند و بهای آن، احتمال کوچکی برای از دست دادن نزدیک‌ترین همسایه‌ی واقعی است.

در این پروژه دو پایگاه برداری به کار رفت: \lr{ChromaDB} برای اجرای محلی و آزمون‌ها، و \lr{Qdrant Cloud} برای سایت. \lr{Qdrant} امکان نگه داشتن بردار متراکم (معنایی) و بردار تنک (کلیدواژه‌ای) در کنار هم را دارد. بردار تنک برداری است با ابعاد برابر اندازه‌ی واژگان که فقط مؤلفه‌های مربوط به واژه‌های موجود در متن آن غیرصفرند.

\section{بازیابی کلیدواژه‌ای \lr{BM25} و ادغام رتبه}
\subsection{روش \lr{BM25}}
جست‌وجوی معنایی در یافتن مترادف‌ها قوی است، ولی نمادهای دقیق مانند $\gamma_{\sub{air}}$، $R(50)$ یا فرمول شیمیایی $\mathrm{LiB_3O_5}$ را خوب تشخیص نمی‌دهد. دلیلش این است که رمزگذار برداری روی متن عمومی آموزش دیده و برای نمادهای کم‌تکرار نمایش دقیقی ندارد. روش کلاسیک \lr{BM25} \cite{bm25} به هر سند بر پایه‌ی تکرار واژه‌های پرسش در آن امتیاز می‌دهد:
\begin{equation}\label{eq:bm25}
\mathrm{BM25}(q,c)=\sum_{t\in q}\mathrm{IDF}(t)\,\frac{f(t,c)\,(k_1+1)}{f(t,c)+k_1\left(1-b+b\,\dfrac{|c|}{\overline{|c|}}\right)} .
\end{equation}
در این رابطه $f(t,c)$ تعداد تکرار واژه‌ی $t$ در بخش $c$، $|c|$ طول بخش، $\overline{|c|}$ میانگین طول بخش‌ها و $k_1$ و $b$ پارامترهای تنظیم‌اند. وزن وارون فراوانی سند،
\begin{equation}\label{eq:idf}
\mathrm{IDF}(t)=\ln\!\left(1+\frac{N-n_t+0.5}{n_t+0.5}\right),
\end{equation}
به واژه‌هایی که در تعداد کمی از $N$ بخش آمده‌اند ($n_t$ کوچک) وزن بیشتری می‌دهد. پس نماد کمیابی مانند $\mathrm{LiB_3O_5}$ وزن بسیار بیشتری از واژه‌ی رایجی مانند \lr{crystal} می‌گیرد. پارامتر $k_1$ اشباع اثر تکرار را کنترل می‌کند (تکرار دهم یک واژه امتیاز کمتری از تکرار اول می‌افزاید) و $b$ اثر طول بخش را بهنجار می‌کند.

\subsection{ادغام معکوس رتبه}
برای ترکیب دو فهرست رتبه‌بندی‌شده، روش ادغام معکوس رتبه (\lr{Reciprocal Rank Fusion} یا \lr{RRF}) \cite{rrf} به کار رفت:
\begin{equation}\label{eq:rrf}
\mathrm{RRF}(c)=\sum_{r\in R}\frac{1}{k+\mathrm{rank}_r(c)},\qquad k=60,
\end{equation}
که در آن $R$ مجموعه‌ی فهرست‌های رتبه‌بندی (اینجا برداری و \lr{BM25}) و $\mathrm{rank}_r(c)$ رتبه‌ی بخش $c$ در فهرست $r$ است. مزیت \lr{RRF} این است که به مقیاس امتیازهای دو روش وابسته نیست و فقط از رتبه‌ها استفاده می‌کند. امتیاز شباهت کسینوسی در بازه‌ی صفر تا یک است، اما امتیاز \lr{BM25} حد بالایی ندارد و ترکیب مستقیم آن‌ها به بهنجارسازی دلخواه نیاز دارد. ثابت $k=60$ اثر تفاوت رتبه‌های بالا را نرم می‌کند، به‌طوری که بخشی که در هر دو فهرست رتبه‌ی متوسطی دارد، می‌تواند از بخشی که فقط در یک فهرست اول است جلو بزند.

\subsection{رتبه‌بندی دوم}
در رتبه‌بندی دوم، تعداد محدودی از نامزدهای اول (مثلاً ۳۰ بخش) با مدلی دقیق‌تر و کندتر دوباره مرتب می‌شوند \cite{nogueira}. برخلاف رمزگذار دوگانه، رتبه‌بند «متقاطع» (\lr{cross-encoder}) پرسش و بخش را با هم می‌خواند و برای هر جفت یک امتیاز مرتبط بودن می‌دهد. این کار دقیق‌تر است، چون مدل می‌تواند واژه‌های پرسش و بخش را مستقیماً با هم مقایسه کند، اما چون برای هر جفت یک اجرای کامل مدل لازم است، فقط روی تعداد کمی نامزد عملی است. یک مدل زبانی بزرگ هم می‌تواند نقش «داور» را بازی کند و برای هر بخش تعیین کند که آیا به پرسش پاسخ می‌دهد یا نه.

\section{مبانی فیزیکی داده‌های هدف}
\subsection{پارامترهای خطوط طیفی}
هر خط جذبی مولکولی در پایگاه \lr{HITRAN} با چند پارامتر توصیف می‌شود \cite{hitran}: مکان خط $\nu$ (بر حسب $\mathrm{cm^{-1}}$)، شدت خط $S$، ضریب پهن‌شدگی فشاری با هوا $\gamma_{\sub{air}}$ و پهن‌شدگی خودی $\gamma_{\sub{self}}$، نمای وابستگی دمایی $n_{\sub{air}}$ و جابه‌جایی فشاری $\delta_{\sub{air}}$. نیم‌پهنای لورنتزی خط در فشار $p$ و دمای $T$ از این پارامترها ساخته می‌شود:
\begin{equation}\label{eq:gamma}
\gamma(p,T)=\left(\frac{T_{\sub{ref}}}{T}\right)^{n_{\sub{air}}}\Big[\gamma_{\sub{air}}\,(p-p_{\sub{self}})+\gamma_{\sub{self}}\,p_{\sub{self}}\Big],
\end{equation}
که $T_{\sub{ref}}=296\ \mathrm{K}$ دمای مرجع \lr{HITRAN} و $p_{\sub{self}}$ فشار جزئی خود گاز است. مکان خط نیز با فشار به‌اندازه‌ی $\delta_{\sub{air}}\,p$ جابه‌جا می‌شود. شکل خط با پروفایل \lr{Voigt} مدل می‌شود که هم‌آمیخت پروفایل گاوسی پهن‌شدگی دوپلر و پروفایل لورنتزی پهن‌شدگی برخوردی است:
\begin{equation}\label{eq:voigt}
V(\nu)=\int_{-\infty}^{\infty}G(\nu';\alpha_D)\,L(\nu-\nu';\gamma)\,d\nu' .
\end{equation}
در فشارهای جوّی نزدیک سطح زمین، سهم لورنتزی غالب است و دقت $\gamma_{\sub{air}}$ مستقیماً دقت مدل‌سازی جذب جوّی را تعیین می‌کند.

نکته‌ی کلیدی برای این پروژه آن است که \lr{HITRAN} برای هر پارامتر، مقاله‌ی منبع را ذکر می‌کند. پس می‌توان پارامتر را با \lr{RAG} از همان مقاله استخراج و با عدد \lr{HITRAN} مقایسه کرد. این ویژگی یک «کلید پاسخ» مستقل و رایگان فراهم می‌کند.

\subsection{ضریب شکست و پاشندگی}
وابستگی ضریب شکست به طول موج در ناحیه‌ی شفافیت معمولاً با فرمول تجربی سلمایر بیان می‌شود:
\begin{equation}\label{eq:sellmeier}
n^{2}(\lambda)=A+\sum_{i}\frac{B_i\lambda^{2}}{\lambda^{2}-C_i},
\end{equation}
که در آن $C_i$ به مربع طول موج‌های تشدید مربوط است. ریشه‌ی فیزیکی این فرمول مدل نوسانگر لورنتز است: هر جمله‌ی کسری سهم یک تشدید الکترونی (در فرابنفش) یا شبکه‌ای (در فروسرخ) را در قطبش‌پذیری محیط نشان می‌دهد. چون این فرمول برازش تجربی است، بیرون از بازه‌ی طول موجی که اندازه‌گیری روی آن انجام شده قابل‌اتکا نیست؛ این بازه بخشی ضروری از هر رکورد داده است.

پایگاه \lr{refractiveindex.info} نه شکل مختلف فرمول پاشندگی (سلمایر، کوشی، چندجمله‌ای و جز آن) را تعریف می‌کند \cite{polyanskiy} که همه در این پروژه پیاده‌سازی شدند. وابستگی دمایی ضریب شکست با ضریب ترمواپتیک $dn/dT$ یا با ضرایب دمایی در خود فرمول بیان می‌شود. از روی $n(\lambda)$ ضریب گروهی و پاشندگی سرعت گروه (\lr{GVD}) هم محاسبه می‌شوند:
\begin{equation}\label{eq:group}
n_g=n-\lambda\frac{dn}{d\lambda},\qquad
\beta_2=\frac{\lambda^{3}}{2\pi c^{2}}\,\frac{d^{2}n}{d\lambda^{2}} .
\end{equation}
ضریب گروهی سرعت حرکت پوش تپ را تعیین می‌کند و $\beta_2$ پهن‌شدن زمانی تپ‌های کوتاه را؛ هر دو در طراحی تبدیل فرکانس با لیزرهای تپی اهمیت دارند.

در بلورهای ناهمسانگرد، ضریب شکست به جهت انتشار و قطبش نور بستگی دارد و با «بیضی‌گون ضریب شکست» توصیف می‌شود \cite{yariv}:
\begin{equation}\label{eq:ellipsoid}
\frac{x^{2}}{n_x^{2}}+\frac{y^{2}}{n_y^{2}}+\frac{z^{2}}{n_z^{2}}=1 .
\end{equation}
در بلور تک‌محوری دو ضریب اصلی برابرند ($n_x=n_y=n_o$، $n_z=n_e$) و در بلور دومحوری هر سه متفاوت‌اند. در بلور تک‌محوری، موج «عادی» با ضریب $n_o$ مستقل از جهت انتشار منتشر می‌شود و موج «غیرعادی» ضریبی دارد که به زاویه‌ی $\theta$ میان بردار موج و محور نوری وابسته است. همین وابستگی زاویه‌ای ابزار اصلی تطبیق فاز دوشکستی است (بخش \ref{sec:pm}). بنابراین، برای هر بلور ناهمسانگرد دست‌کم دو فرمول پاشندگی لازم است و نسبت دادن درست هر فرمول به محور درست، بخشی از کار ساختن پایگاه داده بود.

\subsection{اپتیک غیرخطی مرتبه‌ی دوم و تطبیق فاز}\label{sec:pm}
در محیط غیرمتقارن مرکزی، قطبش القایی جمله‌ی مرتبه‌ی دوم دارد \cite{boyd}:
\begin{equation}\label{eq:pnl}
P_i(\omega_3)=2\varepsilon_0\sum_{jk}d_{ijk}\,E_j(\omega_1)\,E_k(\omega_2).
\end{equation}
تانسور $d_{ij}$ (در نمادگذاری کوتاه‌شده، که جفت اندیس $jk$ با یک اندیس از ۱ تا ۶ جایگزین می‌شود) به گروه نقطه‌ای بلور وابسته است. تقارن بلور بسیاری از مؤلفه‌ها را صفر یا با هم برابر می‌کند؛ برای نمونه، در گروه نقطه‌ای $3m$ (بلورهای \lr{BBO} و لیتیم نیوبات) فقط چند مؤلفه‌ی مستقل باقی می‌ماند. ضریب مؤثر $d_{\sub{eff}}$ از تصویر تانسور روی بردارهای یکه‌ی قطبش سه موج به دست می‌آید:
\begin{equation}\label{eq:deff}
d_{\sub{eff}}=\sum_{ijk}e_{3i}\,d_{ijk}\,e_{1j}\,e_{2k}.
\end{equation}

فرایندهایی مانند تولید هماهنگ دوم (\lr{SHG})، جمع و تفاضل فرکانس (\lr{SFG} و \lr{DFG}) و نوسانگر پارامتری (\lr{OPO}) فقط وقتی کارآمدند که شرط تطبیق فاز برقرار باشد:
\begin{equation}\label{eq:dk}
\Delta k=k_3-k_1-k_2=\frac{2\pi n_3}{\lambda_3}-\frac{2\pi n_1}{\lambda_1}-\frac{2\pi n_2}{\lambda_2}=0 .
\end{equation}
معنای فیزیکی این شرط پایستگی تکانه‌ی فوتون‌هاست: موج قطبش غیرخطی و موج تولیدشده باید هم‌گام حرکت کنند تا سهم لایه‌های پیاپی بلور به‌طور سازنده جمع شود. به دلیل پاشندگی عادی، در یک محیط همسانگرد معمولاً $n_3>n_1,n_2$ است و این شرط خودبه‌خود برقرار نمی‌شود.

در بلور دوشکستی، این شرط با انتخاب زاویه‌ی انتشار (تطبیق فاز بحرانی) یا با تنظیم دما در زاویه‌ی ثابت (تطبیق فاز غیربحرانی، \lr{NCPM}) برقرار می‌شود. در برهم‌کنش «نوع اول» دو موج ورودی قطبش یکسان دارند و در «نوع دوم» قطبش‌هایشان عمود بر هم است. در تطبیق فاز شبه‌فازی (\lr{QPM}) \cite{armstrong}، قطبش بلور به‌صورت دوره‌ای وارون می‌شود (مانند \lr{PPLN} و \lr{PPKTP}) و دوره‌ی لازم و ضریب مؤثر از رابطه‌ی زیر به دست می‌آیند:
\begin{equation}\label{eq:qpm}
\Lambda=\frac{2\pi}{|\Delta k|},\qquad d_{\sub{eff}}^{\sub{QPM}}=\frac{2}{\pi}\,|d| .
\end{equation}
ضریب $2/\pi$ از نخستین جمله‌ی سری فوریه‌ی تابع مربعی وارون‌شدن قطبش می‌آید. مزیت بزرگ \lr{QPM} این است که می‌توان هر سه موج را با یک قطبش انتخاب کرد و از بزرگ‌ترین مؤلفه‌ی تانسور (مثلاً $d_{33}$) بهره برد، که در تطبیق فاز دوشکستی در دسترس نیست.

کمیت‌های عملی دیگر عبارت‌اند از جدایی مسیر انرژی (\lr{walk-off})، و پهنای پذیرش زاویه‌ای، طیفی و دمایی (بخش \ref{sec:acc}). همه‌ی این محاسبه‌ها به $n(\lambda,T)$ دقیق وابسته‌اند. خطای $10^{-3}$ در ضریب شکست می‌تواند زاویه‌ی تطبیق فاز را به‌اندازه‌ی محسوسی جابه‌جا کند. به همین دلیل، داده‌ی ورودی باید با منبع آزموده شده باشد.

\subsection{بازده تبدیل، طول همدوسی و پهنای پذیرش}\label{sec:acc}
در تقریب موج تخت و دمش تهی‌نشده، شدت هماهنگ دوم پس از طول $L$ از بلور از رابطه‌ی زیر پیروی می‌کند \cite{boyd}:
\begin{equation}\label{eq:shg}
I_{2\omega}\propto\frac{d_{\sub{eff}}^{2}\,L^{2}}{n_{\omega}^{2}\,n_{2\omega}\,\lambda_{\omega}^{2}}\;I_{\omega}^{2}\;\sinc^{2}\!\left(\frac{\Delta k L}{2}\right).
\end{equation}
این رابطه سه نتیجه‌ی عملی دارد. نخست، بازده با $d_{\sub{eff}}^{2}$ متناسب است. پس خطای ۱۰ درصدی در ضریب غیرخطی، خطای حدود ۲۰ درصدی در توان پیش‌بینی‌شده می‌دهد. دوم، اگر $\Delta k\neq0$ باشد، توان هماهنگ دوم با طول بلور به‌طور نوسانی تغییر می‌کند و بیشینه‌ی آن در «طول همدوسی» رخ می‌دهد:
\begin{equation}\label{eq:lc}
L_c=\frac{\pi}{|\Delta k|}.
\end{equation}
سوم، چون تابع $\sinc^{2}$ در $|\Delta k L/2|\approx1.39$ به نصف بیشینه می‌رسد، هر انحراف از شرایط تطبیق فاز (زاویه، طول موج یا دما) فقط تا حد معینی تحمل می‌شود. این حد «پهنای پذیرش» نام دارد و از مشتق $\Delta k$ نسبت به آن متغیر به دست می‌آید. برای نمونه، پهنای پذیرش زاویه‌ای (پهنای کامل در نصف بیشینه) برابر است با:
\begin{equation}\label{eq:acc}
\Delta\theta\cdot L\simeq\frac{2\times2.78}{|\partial\Delta k/\partial\theta|}.
\end{equation}
پهنای پذیرش دمایی و طیفی هم به همین شکل از $\partial\Delta k/\partial T$ و $\partial\Delta k/\partial\lambda$ محاسبه می‌شوند. چون این مشتق‌ها مستقیماً از مشتق‌های $n(\lambda,T)$ ساخته می‌شوند، خطای کوچک در شیب فرمول پاشندگی، حتی اگر خود مقدار $n$ درست باشد، پهنای پذیرش را جابه‌جا می‌کند. این روابط نشان می‌دهند که چرا یک فرمول پاشندگی نادرست یا یک مدل ترمواپتیک ناقص، مستقیماً به خطا در طراحی آزمایش منجر می‌شود. رابطه‌ی (\ref{eq:shg}) برای باریکه‌ی کانونی‌شده باید با نظریه‌ی \lr{Boyd-Kleinman} \cite{boydkleinman} اصلاح شود که در این پروژه پیاده‌سازی نشده است (بخش \ref{sec:limits}). در فصل ۵ این روابط برای بلور \lr{BBO} به‌طور عددی حل شده‌اند.

\subsection{فروتبدیل پارامتری خودبه‌خودی}
در فروتبدیل پارامتری خودبه‌خودی (\lr{SPDC}) یک فوتون دمش با بسامد $\omega_p$ در بلور غیرخطی به دو فوتون «علامت» و «بیکار» تبدیل می‌شود. پایستگی انرژی و تکانه می‌دهد:
\begin{equation}\label{eq:spdc}
\omega_p=\omega_s+\omega_i,\qquad \Delta k=k_p-k_s-k_i-\frac{2\pi}{\Lambda}\approx0 ,
\end{equation}
که جمله‌ی $2\pi/\Lambda$ فقط در بلورهای قطبیده‌ی دوره‌ای حضور دارد. حالت دوفوتونی با «دامنه‌ی طیفی مشترک» توصیف می‌شود که حاصل‌ضرب پوش طیفی دمش و تابع تطبیق فاز است \cite{bennink}:
\begin{equation}\label{eq:jsa}
f(\omega_s,\omega_i)=\alpha(\omega_s+\omega_i)\;\phi(\omega_s,\omega_i),\qquad
\phi\propto\sinc\!\left(\frac{\Delta k L}{2}\right).
\end{equation}
با تجزیه‌ی اشمیت این تابع، $f=\sum_j\sqrt{\lambda_j}\,u_j(\omega_s)\,v_j(\omega_i)$، عدد اشمیت $K=1/\sum_j\lambda_j^{2}$ و خلوص طیفی تک‌فوتون $P=1/K$ به دست می‌آیند. $K=1$ یعنی دو فوتون از نظر طیفی هیچ همبستگی ندارند، که برای منابع تک‌فوتونِ اعلام‌شده مطلوب است. این کمیت‌ها نیز کاملاً به $n(\lambda,T)$ و دوره‌ی $\Lambda$ وابسته‌اند.

\section{پیشینه‌ی پژوهش}
\textbf{استخراج خودکار داده از متن علمی.} ابزار \lr{ChemDataExtractor} \cite{cde} با قواعد زبانی و یادگیری ماشین، پایگاه‌های داده‌ای با ده‌ها هزار رکورد از خواص مواد ساخته است. این روش در مقیاس قوی است، اما برای فرمول‌های چندضریبی و جدول‌های پیچیده دقت محدودی دارد، زیرا قواعد آن برای الگوهای «نام ماده، خاصیت، مقدار» طراحی شده‌اند و نه برای مجموعه‌ای از ضرایب که فقط با هم معنا دارند.

\textbf{پایگاه‌های مواد.} \lr{Materials Project} \cite{mp} داده‌ی محاسباتی (\lr{DFT}) برای بیش از صد هزار ماده دارد. \lr{refractiveindex.info} \cite{polyanskiy} مرجع آزاد ضریب شکست است و کتاب‌های دستی \lr{Dmitriev} و همکاران \cite{dmitriev} و \lr{Nikogosyan} \cite{nikogosyan} مرجع کلاسیک بلورهای غیرخطی‌اند. این مراجع داده را گردآوری کرده‌اند، اما ارتباط مستقیمی با متن کامل مقاله‌ی منبع و ابزار محاسبه ندارند.

\textbf{بازیابی متراکم و رتبه‌بندی.} نشان داده شده است که رمزگذارهای دوگانه‌ی آموزش‌دیده با یادگیری تقابلی در پرسش و پاسخ با دامنه‌ی باز از \lr{BM25} بهتر عمل می‌کنند \cite{dpr}، و رتبه‌بند متقاطع مبتنی بر ترنسفورمر کیفیت رتبه‌ی اول را بهتر می‌کند \cite{nogueira}. با این حال، این نتایج بیشتر روی متن عمومی (مانند ویکی‌پدیا) به دست آمده‌اند و برای متن علمی پر از نماد نیازمند آزمون مستقل‌اند؛ پرسش دوم این پژوهش از همین‌جا می‌آید.

\textbf{دستیارهای پژوهشی مبتنی بر \lr{RAG}.} سامانه‌هایی مانند \lr{PaperQA2} \cite{paperqa} با بازیابی و خواندن تکراری مقاله‌ها پاسخ مستند می‌نویسند و در برخی بنچمارک‌ها با متخصصان انسانی مقایسه شده‌اند. روش‌های ارزیابی \lr{RAG} مانند \lr{RAGAS} \cite{ragas} هم برای سنجش وفاداری پاسخ به منبع پیشنهاد شده‌اند.

کار حاضر در محل تلاقی این دسته‌ها است: بازیابی متن مقاله، داده‌ی عددی آزموده و ابزار محاسبه. هیچ‌کدام از این دسته‌ها به‌تنهایی این زنجیره را ندارد، ولی هر کدام در حوزه‌ی خودش از این پروژه قوی‌تر و جاافتاده‌تر است.

% ======================================== فصل ۳ ======================================================
\chapter{روش و طراحی سامانه}

\section{معماری کلی}
سامانه سه لایه دارد. \textbf{لایه‌ی دانش} مقاله‌ها را به بخش‌های قابل‌جست‌وجو تبدیل و در پایگاه داده‌های برداری نمایه می‌کند. \textbf{لایه‌ی پرسش و استخراج} بخش‌های مرتبط را بازیابی می‌کند و پاسخ مستند یا رکورد ساخت‌یافته می‌سازد. \textbf{لایه‌ی داده و ابزار} عددهای آزموده را در پایگاه داده‌ی فیزیکی نگه می‌دارد و به محاسبه‌گرها می‌دهد. شکل \ref{fig:arch} جریان داده را نشان می‌دهد. جدا بودن این سه لایه یک تصمیم طراحی آگاهانه بود: لایه‌ی دانش را می‌توان بدون تغییر لایه‌های دیگر با مجموعه‌ی تازه‌ای از مقاله‌ها پر کرد، و محاسبه‌گرها فقط از داده‌ی آزموده تغذیه می‌شوند، نه از خروجی مستقیم مدل زبانی.

\begin{figure}[htbp]
\centering
\archfigure
\caption{معماری سامانه در سه مسیر نمایه‌سازی، پرسش و پاسخ، و داده و ابزار.}
\label{fig:arch}
\end{figure}

\section{مجموعه‌ی اسناد}
\subsection{مجموعه‌ی آزمایشی \lr{HITRAN}}
ایده‌ی ارزیابی این مرحله از خود \lr{HITRAN} گرفته شد. این پایگاه از مقاله‌ها گردآوری شده و برای هر پارامتر مقاله‌ی منبع را ذکر می‌کند. پس پارامترهای استخراج‌شده با \lr{RAG} از همان مقاله‌ها را می‌توان با کتابخانه‌ی رسمی \lr{HAPI} \cite{hapi} با عدد \lr{HITRAN} مقایسه کرد، بدون نیاز به برچسب‌گذاری دستی گسترده. مجموعه‌ی آزمایشی شامل مقاله‌های دسترسی آزاد \lr{arXiv} درباره‌ی پهن‌شدگی و شدت خطوط $\mathrm{CO_2}$ و $\mathrm{CH_4}$ بود و ۱۶ پرسش مرجع با صفحه‌ی پاسخ و مقدار درست برای آن نوشته شد. یکی از این پرسش‌ها عمداً بی‌جواب طراحی شد تا رفتار سامانه در نبود شاهد هم سنجیده شود.

\subsection{مجموعه‌ی بلورهای نوری}
مجموعه‌ی اصلی پروژه حوزه‌ی اپتیک خطی و غیرخطی است. فهرست اولیه از مقاله‌های منبع \lr{refractiveindex.info}، کتاب‌های دستی اپتیک غیرخطی و جست‌وجوی هدفمند برای هر خانواده‌ی بلور ساخته شد. مشخصات مقاله‌ها (\lr{DOI}، سال، نویسندگان، مجله و چکیده) از \lr{Crossref} و \lr{OpenAlex} گرفته شد. جدول \ref{tab:corpus} ترکیب مجموعه را نشان می‌دهد.

\begin{table}[htbp]
\centering
\caption{ترکیب مجموعه‌ی بلورهای نوری (شمارش مستقیم از پایگاه \lr{Qdrant}، ۳ آبان ۱۴۰۵).}
\label{tab:corpus}
\renewcommand{\arraystretch}{1.3}
\resizebox{\textwidth}{!}{%
\begin{tabular}{lccc}
\toprule
نوع سند & تعداد & سهم & بخش نمایه‌شده \\
\midrule
متن کامل \lr{PDF} & ۲۲۱۴ & ۷۷/۸ درصد & ۶۷۵۵۹ (۱۱۴۲۱ جدول) \\
فقط چکیده (\lr{PDF} آزاد یا قانونی در دسترس نبود) & ۵۹۸ & ۲۱/۰ درصد & ۵۹۸ \\
فقط مشخصات کتاب‌شناختی & ۳۳ & ۱/۲ درصد & ۳۳ \\
\midrule
جمع & ۲۸۴۵ & ۱۰۰ درصد & ۶۸۱۹۰ \\
\bottomrule
\end{tabular}}
\end{table}

از ۲۷۲۸ مقاله‌ی مجله‌ای، ۱۳۳۳ مقاله (۴۸/۹ درصد) در مجله‌های چارک \lr{Q1} و ۹۳۸ مقاله (۳۴/۴ درصد) در چارک \lr{Q2} فهرست \lr{SCImago} چاپ شده‌اند. بیشتر ۲۳۲ مقاله‌ای که در \lr{SCImago} یافت نشدند، با نام قدیمی مجله‌های معتبر (مانند \lr{Physical Review B, Condensed Matter}) ثبت شده‌اند. حدود ۱۷ درصد مجموعه (۴۸۱ سند) پیش از ۱۹۹۰ منتشر شده است. این مقاله‌ها اغلب اسکن تصویری‌اند و به \lr{OCR} نیاز دارند. وجود این مقاله‌های قدیمی اجتناب‌ناپذیر است، چون بسیاری از اندازه‌گیری‌های پایه‌ی ضریب شکست و ضرایب غیرخطی در همان دهه‌ها انجام شده‌اند و هنوز مرجع‌اند.

مجموعه در ۹ قفسه‌ی موضوعی مرتب شده است، از جمله بورات‌ها (۵۴۱ سند)، کالکوپیریت‌ها و کالکوژنیدهای فروسرخ میانی (۵۰۳)، نظریه و مدل‌های پاشندگی (۴۱۶) و نیوبات‌ها و تانتالات‌ها (۳۳۹).

یک تصمیم روش‌شناختی مهم این بود که داده‌ی \lr{refractiveindex.info} \textbf{عمداً وارد جست‌وجو نشود}. این داده کلید ارزیابی است. اگر وارد نمایه می‌شد، سامانه پاسخ را از خود کلید پیدا می‌کرد و ارزیابی بی‌معنا می‌شد. این همان اصل جدا نگه داشتن داده‌ی آزمون از داده‌ی سامانه است که در هر ارزیابی تجربی رعایت می‌شود.

\section{تبدیل سند و بخش‌بندی}
\lr{PDF}‌ها با کتابخانه‌ی \lr{Docling} \cite{docling} به متن ساختاریافته تبدیل می‌شوند: عنوان بخش، پاراگراف، جدول با سطر و ستون، و کادر محدودکننده‌ی هر عنصر روی صفحه. قالب \lr{PDF} برای نمایش طراحی شده است، نه برای استخراج متن: در آن فقط جای هر نویسه روی صفحه ثبت می‌شود و ترتیب خواندن، مرز ستون‌ها و ساختار جدول باید از روی چیدمان بازسازی شوند. چند مشکل ویژه‌ی متن فیزیکی در این مرحله حل شد:
\begin{itemize}
\item \textbf{حروف یونانی گمشده:} \lr{PDF}‌های ساخته‌شده با \lr{Word} حروف یونانی را با فونت \lr{Symbol} و در ناحیه‌ی \lr{Private Use} یونیکد ذخیره می‌کنند (مثلاً $\nu$ به‌صورت \lr{U+F06E}). این نویسه‌ها برای مدل زبانی و جست‌وجو نامرئی‌اند. با جدول استاندارد \lr{Adobe Symbol Encoding} به یونیکد درست برگردانده شدند. اهمیت این ترمیم در متن طیف‌سنجی آشکار است، جایی که $\nu_3$ و $\gamma$ و $\delta$ نام کمیت‌های اصلی‌اند.
\item \textbf{زیرنویس فرمول شیمیایی:} $\mathrm{LiB_3O_5}$ که پس از استخراج به‌صورت «\lr{LiB 3 O 5}» درمی‌آید، یکسان‌سازی شد تا جست‌وجوی کلیدواژه‌ای کار کند.
\item \textbf{جدول‌ها:} هر جدول به شکل «سطر، ستون = مقدار» سریالایز می‌شود و سطر عنوان در هر قطعه تکرار می‌شود، تا هر عدد همراه با معنایش بازیابی شود. بدون این کار، قطعه‌ای از وسط یک جدول فقط فهرستی از عددهای بی‌نام بود.
\item \textbf{معادله‌ها:} معادله‌های تصویری ۴۶ مقاله با مدل \lr{CodeFormula} روی \lr{GPU} به \lr{LaTeX} تبدیل شدند.
\item \textbf{\lr{OCR}:} مقاله‌های اسکن‌شده با \lr{OCR} دوباره خوانده شدند. هر مقاله فقط وقتی جایگزین شد که متنش بهتر شده بود و هیچ عددی از جدول‌هایش گم نشده بود (۴۵ مقاله از ۶۱ مقاله‌ی دور اول). کل کتاب \lr{Dmitriev} و همکاران (۴۲۶ صفحه) هم \lr{OCR} شد.
\end{itemize}

بخش‌بندی با \lr{HybridChunker} کتابخانه‌ی \lr{Docling} انجام می‌شود. این روش مرز بخش‌ها را بر پایه‌ی ساختار سند (بخش، پاراگراف، جدول) می‌گذارد و هر بخش را زیر سقف ۵۱۲ توکن مدل برداری نگه می‌دارد. انتخاب اندازه‌ی بخش یک مصالحه است: بخش کوچک دقیق‌تر بازیابی می‌شود ولی ممکن است بافت لازم برای فهم عدد (مثلاً شرایط اندازه‌گیری) را نداشته باشد، و بخش بزرگ بافت کامل‌تری دارد ولی بردارش میانگین چند موضوع است. سقف ۵۱۲ توکن از حداکثر طول ورودی مدل برداری می‌آید و مرزبندی ساختاری از شکستن پاراگراف‌ها در وسط جلوگیری می‌کند.

هر بخش این فراداده را دارد: شناسه‌ی سند، شماره‌ی صفحه، عنوان بخش، نوع (متن یا جدول) و مختصات روی صفحه. همین مختصات بعداً برای رنگی کردن محل منبع روی تصویر صفحه‌ی \lr{PDF} به کار می‌رود، تا کاربر بتواند عدد را در خود مقاله ببیند.

\section{نمایه در پایگاه داده‌های برداری}
هر بخش با مدل \lr{bge-base-en-v1.5} به بردار ۷۶۸ بعدی تبدیل می‌شود. این مدل به این دلیل انتخاب شد که محلی و رایگان اجرا می‌شود (تکرارپذیری و صرفه‌جویی در سهمیه‌ی \lr{API}) و در بنچمارک‌های بازیابی انگلیسی عملکرد خوبی دارد \cite{bge}. در اجرای محلی، بردارها و متن در \lr{ChromaDB} و یک \lr{docstore} در \lr{LlamaIndex} ذخیره می‌شوند و نمایه‌ی \lr{BM25} با کتابخانه‌ی \lr{bm25s} ساخته می‌شود. برای سایت، \textbf{همان بردارهای ذخیره‌شده} به \lr{Qdrant Cloud} منتقل شدند و در کنار هر بردار متراکم یک بردار تنک \lr{BM25} سمت سرور نگه داشته شد. یکسان بودن بردارها تضمین می‌کند که نتیجه‌ی ارزیابی محلی و سایت قابل‌مقایسه باشد. محاسبه‌ی دوباره‌ی بردارها با نسخه‌ی دیگری از مدل یا کتابخانه می‌توانست تفاوت‌های کوچکی ایجاد کند که مقایسه را مخدوش کند.

\section{بازیابی}
\subsection{بازیابی ترکیبی}
برای هر پرسش، جست‌وجوی برداری و \lr{BM25} جداگانه اجرا و نتیجه‌ها با \lr{RRF} (رابطه‌ی (\ref{eq:rrf})) ادغام می‌شوند. در اجرای محلی ۸ بخش برتر به مدل زبانی داده می‌شود. انگیزه‌ی این طراحی، ماهیت دوگانه‌ی متن فیزیکی است: پرسش «پهن‌شدگی با بخار آب» به جست‌وجوی معنایی نیاز دارد، چون مقاله ممکن است عبارت \lr{H2O-broadening} یا \lr{water-vapor broadened} را به کار برده باشد، و پرسش درباره‌ی $\nu_3$ یا $R(50)$ به تطابق دقیق نماد.

\subsection{گسترش نام ماده}
نام کوتاه، نام کامل و نام فارسی بلور به فرمول شیمیایی نگاشته می‌شود. برای نمونه «\lr{BBO}» و «باریم بورات» هر دو به $\mathrm{BaB_2O_4}$ نگاشته می‌شوند. یک جدول واحد نام‌ها برای جست‌وجو و محاسبه‌گرها استفاده می‌شود. در این جدول، طولانی‌ترین نام اول تطبیق داده می‌شود تا «\lr{cesium lithium borate}» با \lr{LBO} یا «\lr{alpha-BBO}» با \lr{BBO} اشتباه نشود. فرمول‌ها با حروف بزرگ و کوچک خودشان تطبیق داده می‌شوند تا «\lr{gap}» با \lr{GaP} یکی نشود. اگر پرسش نام ماده داشته باشد، یک جست‌وجوی جداگانه روی همان ماده هم اجرا می‌شود. این جست‌وجوی جداگانه برای پرسش‌هایی لازم است که بیشتر واژه‌هایشان عمومی‌اند (مانند «ضریب شکست») و نام ماده فقط یک واژه از آن‌هاست.

\subsection{پرسش فارسی}
مقاله‌ها انگلیسی‌اند و مدل برداری هم انگلیسی است. پس پرسش فارسی اول با یک مدل زبانی به انگلیسی ترجمه می‌شود. ترجمه‌ی تحت‌اللفظی اصطلاحات فیزیک اغلب غلط است. برای نمونه «عدد موج» \lr{wavenumber} است و نه \lr{wavelength}، و «غیرعادی» در اپتیک بلور \lr{extraordinary} است و نه \lr{abnormal}. برای همین یک واژه‌نامه‌ی ثابت فیزیکی به مدل ترجمه داده می‌شود. در اجرای محلی ترجمه با همان مدل محلی (\lr{Qwen3-4B}) و در سایت با \lr{Llama 3.3 70B} روی \lr{Workers AI} انجام می‌شود. ترجمه و بردار پرسش‌ها در حافظه‌ی نهان نگه داشته می‌شوند تا پرسش تکراری هزینه‌ی دوباره نداشته باشد.

\subsection{رتبه‌بندی دوم}
در سایت، ۳۰ نامزد اول می‌توانند با یک مدل داور (\lr{SubQ}) دوباره مرتب شوند. این مدل برای هر بخش تشخیص می‌دهد که آیا بخش به پرسش پاسخ می‌دهد یا نه. رتبه‌بند متقاطع \lr{bge-reranker-base} هم آزموده شد (بخش \ref{sec:ablation}).

\subsection{نگهبان «پیدا نشد»}
اگر نام ماده یا واژه‌های اصلی پرسش در هیچ نتیجه‌ای نباشد، یا شباهت برداری بهترین نتیجه از یک آستانه کمتر باشد، سامانه صریحاً می‌گوید منبعی پیدا نشد. این نگهبان در برابر خطای رایج سامانه‌های بازیابی طراحی شده است: جست‌وجوی برداری همیشه نزدیک‌ترین بردارها را برمی‌گرداند، حتی اگر هیچ‌کدام واقعاً مرتبط نباشند.

\section{تولید پاسخ و استخراج ساخت‌یافته}
\subsection{پاسخ با ارجاع}
مدل زبانی (\lr{Qwen3-4B} محلی از راه \lr{Ollama} \cite{qwen3}، یا \lr{Gemini}) پرسش و بخش‌های بازیابی‌شده را می‌بیند. دستور مدل این است که فقط از همین بخش‌ها پاسخ دهد، برای هر عدد منبع را به شکل «نویسنده (سال)، صفحه» بیاورد و اگر شاهدی نیست بگوید «پیدا نشد». برای مقایسه، همان مدل بدون منبع هم پاسخ می‌دهد. خروجی دیگری به نام «بسته‌ی منابع» هم ساخته می‌شود: فایلی شامل پرسش و بخش‌های منبع که می‌توان آن را در هر چت هوش مصنوعی بارگذاری کرد، بدون مصرف \lr{API}.

قواعد اصلی دستور (\lr{prompt}) مدل، که عیناً در کد پروژه به کار رفته‌اند، چنین‌اند:
\begin{enumerate}
\item فقط از منابع شماره‌دار داده‌شده پاسخ بده و هر جمله را با شماره‌ی منبعش ارجاع بده.
\item مقادیر عددی را دقیقاً همان‌طور که نوشته شده‌اند، همراه با واحد، عدم‌قطعیت و شرایط اندازه‌گیری (دما، فشار، گاز پهن‌کننده، باند طیفی) گزارش کن.
\item اگر منبع مقدار تجربی را از مقدار نظری یا محاسبه‌شده جدا کرده، این تمایز را حفظ کن.
\item اگر منابع پاسخ را ندارند، دقیقاً بنویس: «\lr{Not found in the provided sources}».
\item ارجاع را با شماره‌ی صفحه بنویس، مانند \lr{(Source 2, p. 5)}، و پاسخ را به فارسی بنویس، ولی نمادها، واحدها و اعداد را همان‌طور که در منبع آمده نگه دار.
\end{enumerate}
قاعده‌ی دوم برای داده‌ی فیزیکی اهمیت ویژه دارد. مقدار «\lr{0.0712(3)}» بدون عدم‌قطعیت، یا ضریب پهن‌شدگی بدون نام گاز پهن‌کننده، از نظر فیزیکی ناقص است. قاعده‌ی سوم هم از آن رو لازم است که بسیاری از مقاله‌های طیف‌سنجی مقدار اندازه‌گیری‌شده را کنار مقدار محاسبه‌شده‌ی نظری در یک جدول می‌آورند و جابه‌جایی این دو خطایی جدی است.

\subsection{استخراج ساخت‌یافته}
برای استخراج خواص، مدل زبانی یک قالب داده‌ی \lr{Pydantic} را پر می‌کند. فیلدهای اصلی این قالب عبارت‌اند از: ماده، ایزوتوپولوگ، نام خاصیت، مقدار به همان شکلی که نوشته شده، مقدار عددی، عدم‌قطعیت، واحد، موقعیت طیفی، دما، فشار، گاز پهن‌کننده، روش (تجربی، نظری یا گردآوری‌شده) و یک نقل‌قول عینی از متن به‌عنوان شاهد. نگه داشتن «مقدار به همان شکلی که نوشته شده» در کنار مقدار عددی، امکان بررسی بعدی را فراهم می‌کند. سپس خط لوله شناسه‌ی سند، بخش و صفحه را اضافه می‌کند و سه بررسی خودکار انجام می‌دهد (جدول \ref{tab:checks}).

\begin{table}[htbp]
\centering
\caption{بررسی‌های خودکار روی هر مقدار استخراج‌شده.}
\label{tab:checks}
\renewcommand{\arraystretch}{1.3}
\small
\begin{tabular}{p{3.3cm}p{6.0cm}p{5.3cm}}
\toprule
بررسی & پرسش & خطایی که می‌گیرد \\
\midrule
\lr{value\_in\_source} & آیا عدد، همان‌طور که نوشته شده، در متن بخش منبع هست؟ & عددسازی مدل \\
\lr{evidence\_verified} & آیا جمله‌ی شاهد عیناً در متن منبع هست؟ & شاهد ساختگی \\
\lr{plausible} & آیا مقدار در بازه‌ی فیزیکی معقول آن خاصیت است؟ & خطای واحد، برداشت اشتباه (مثلاً «تعداد خط‌ها» به‌جای ضریب) \\
\bottomrule
\end{tabular}
\end{table}

بازه‌های فیزیکی برای هر مجموعه در یک فایل پیکربندی تعریف شده‌اند. برای نمونه، ضریب شکست باید در بازه‌ی $1<n<5$ باشد. این بررسی‌ها مکمل یکدیگرند: بررسی اول جلوی عددی را می‌گیرد که در منبع نیست، بررسی دوم جلوی شاهد ساختگی را، و بررسی سوم عددی را که در منبع هست ولی به خاصیت دیگری تعلق دارد.

\subsection{استخراج فرمول پاشندگی}
بسیاری از مقاله‌های اپتیکی ضریب شکست را فقط به شکل نمودار نشان می‌دهند، اما ضرایب فرمول برازش‌شده را در جدول یا متن می‌آورند. سامانه این ضرایب را استخراج می‌کند و با آن‌ها $n(\lambda)$ را در هر طول موجی حساب می‌کند. هر فرمول استخراج‌شده سه بررسی دارد: همه‌ی ضرایب در متن منبع هستند؟ فرمول از نظر فیزیکی معقول است ($1<n<5$، بدون تشدید در بازه و با $|dn/d\lambda|$ کوچک)؟ و اختلافش با مدخل همان مقاله در \lr{refractiveindex.info} چقدر است؟

\section{پایگاه داده‌ی فیزیکی آزموده}
نتیجه‌ی بخش \ref{sec:negative} نشان داد که استخراج کاملاً خودکار ضرایب به دقت لازم نمی‌رسد. پس برای داده‌ای که به محاسبه‌گرها وارد می‌شود، روش نیمه‌دستی انتخاب شد. سامانه بخش منبع را پیدا می‌کند، عدد خوانده و ثبت می‌شود و سپس با یک «آزمون» درستی‌آزمایی می‌شود. قاعده‌ی کلیدی این است: \textbf{هیچ عددی بدون آزمون با عدد خود مقاله وارد پایگاه نمی‌شود.} هر رکورد شناسه‌ی سند، شماره‌ی صفحه و جای دقیق داده (مثلاً «جدول ۲، ستون سوم») را دارد. هر فرمول هم یک مقدار گزارش‌شده در خود مقاله همراه دارد، مانند زاویه‌ی تطبیق فاز، ضریب شکست در یک طول موج یا دوره‌ی \lr{QPM}. آزمون خودکار بررسی می‌کند که محاسبه با ضرایب ثبت‌شده همان عدد را بازتولید کند. این روش خطای خواندن ضرایب را می‌گیرد، چون یک ضریب نادرست تقریباً همیشه عدد گزارش‌شده را تغییر می‌دهد. جدول \ref{tab:assets} دارایی‌های پایگاه را نشان می‌دهد.

\begin{table}[htbp]
\centering
\caption{دارایی‌های پایگاه داده‌ی فیزیکی (۳ آبان ۱۴۰۵).}
\label{tab:assets}
\renewcommand{\arraystretch}{1.3}
\small
\begin{tabular}{p{3.8cm}p{2.4cm}p{8.4cm}}
\toprule
داده & پوشش & منبع \\
\midrule
منحنی ضریب شکست & ۲۳۳۰ ماده & \lr{refractiveindex.info} (مجوز \lr{CC0}) به‌اضافه‌ی ۹۷ فرمول برای ۸۴ ماده که مستقیماً از مقاله‌های مجموعه خوانده شد \\
تانسور غیرخطی $d_{ij}$ & ۱۰۵ بلور & \lr{Roberts 1992} \cite{roberts}، کتاب‌های \lr{Dmitriev} \cite{dmitriev} و \lr{Nikogosyan} \cite{nikogosyan} و اندازه‌گیری‌های مستقیم \\
مدل ترمواپتیک & ۶۷ بلور (۷۹ مدل) & فرمول‌های دمایی \lr{Kato} و همکاران، \lr{Ghosh} و اندازه‌گیری‌های مستقیم \\
بلور ناهمسانگرد در محاسبه‌گر & ۱۸۶ بلور & هر بلوری که فرمول دو یا سه محور دارد \\
کارت مشخصات عملی (شفافیت، رطوبت‌گیری، آستانه‌ی آسیب) & ۵۸ بلور & \lr{Nikogosyan 2005} و مقاله‌های مجموعه \\
\bottomrule
\end{tabular}
\end{table}

برای نمونه، رکورد بلور \lr{BBO} در فایل \lr{\texttt{nonlinear\_coefficients.yml}} به این شکل است:
\codebox{\texttt{%
- material: BaB2O4\\
\hspace*{1em}point\_group: 3m\\
\hspace*{1em}d: \{22: 2.2, 21: -2.2, 16: -2.2, 31: 0.1, 32: 0.1, 15: 0.1, 24: 0.1\}\\
\hspace*{1em}sources:\\
\hspace*{2em}- \{doc\_id: 10.1109\_3.55534, cite: "Eckardt et al. (1990)", pdf\_page: 10,\\
\hspace*{3em}what: "|d22| = 2.2 pm/V (Table III), from the measured d\_eff = 1.94 pm/V\\
\hspace*{3em}of type-I SHG at 1.064 um"\}\\
\hspace*{2em}- \{doc\_id: 10.1109\_3.159516, cite: "Roberts (1992)", pdf\_page: 13,\\
\hspace*{3em}what: "d31 = 0.1 pm/V and the relative sign d21 = -d22 (Table VI)"\}\\
\hspace*{1em}notes: "3m: d21 = d16 = -d22, d31 = d32 = d15 = d24"}}

این رکورد سه ویژگی طراحی پایگاه را نشان می‌دهد. نخست، همه‌ی مؤلفه‌های غیرصفر تانسور نوشته شده‌اند و مؤلفه‌هایی که تقارن گروه نقطه‌ای $3m$ با هم برابر می‌کند، در یادداشت مشخص شده‌اند. دوم، هر عدد به شناسه‌ی سند (\lr{DOI})، شماره‌ی صفحه‌ی \lr{PDF} و جدول مشخصی ارجاع دارد. سوم، علامت نسبی مؤلفه‌ها از منبع جداگانه‌ای (\lr{Roberts 1992}) آمده و این منبع هم ثبت شده است. سایت $d_{\sub{eff}}$ را مستقیماً از جمع $\sum e_{3i}d_{ijk}e_{1j}e_{2k}$ (رابطه‌ی (\ref{eq:deff})) روی بردارهای میدان سه موج محاسبه می‌کند، پس نیازی به انتخاب دستی فرمول $d_{\sub{eff}}$ برای هر گروه نقطه‌ای نیست. این روش احتمال خطای نقل فرمول‌های $d_{\sub{eff}}$ از جدول‌های کتاب‌ها را هم حذف می‌کند. قاب محورهای بلورهای دومحوری (که در ادبیات گاهی جابه‌جا آمده است) با مقادیر اندازه‌گیری‌شده‌ی $d_{\sub{eff}}$ در مقاله‌های مجموعه تثبیت شده و آزمون‌های خودکار آن را دوباره محاسبه می‌کنند.

% ======================================== فصل ۴ ======================================================
\chapter{پیاده‌سازی}

\section{ساختار نرم‌افزار}
کد پروژه در یک مخزن \lr{Git} نگه داشته می‌شود (\lr{\texttt{github.com/\allowbreak aminsadidi/\allowbreak RAG}}) و از ۵ مهر ۱۴۰۵ تا زمان نگارش این گزارش ۲۹۷ نسخه‌ی ثبت‌شده (\lr{commit}) دارد. نرم‌افزار چهار بخش اصلی دارد که در جدول \ref{tab:software} آمده‌اند.

\begin{table}[htbp]
\centering
\caption{بخش‌های نرم‌افزار.}
\label{tab:software}
\renewcommand{\arraystretch}{1.3}
\small
\begin{tabular}{p{3.7cm}p{3.7cm}p{7.2cm}}
\toprule
بخش & زبان و محل اجرا & کارکرد \\
\midrule
\lr{\texttt{src/\allowbreak matrag/\allowbreak }} & \lr{Python}؛ \lr{Google Colab} یا رایانه‌ی محلی & خط لوله‌ی داده: تبدیل، بخش‌بندی، بردارسازی، نمایه، بازیابی، پاسخ، استخراج، موتور فرمول و ارزیابی \\
\lr{\texttt{notebooks/\allowbreak }} & \lr{Google Colab} & اجرای همان خط لوله روی مجموعه‌ی \lr{Google Drive}؛ چند نوت‌بوک کارگر برای تبدیل موازی \\
\lr{\texttt{web/\allowbreak }} & \lr{JavaScript}؛ \lr{Cloudflare Workers} & سایت: جست‌وجو در \lr{Qdrant}، نمایش صفحه‌ی \lr{PDF} و محاسبه‌گرهایی که در مرورگر اجرا می‌شوند \\
\lr{\texttt{research\_package/\allowbreak }} & \lr{Python} & نسخه‌ی مرجع محاسبه‌های تطبیق فاز و \lr{SPDC} برای آزمودن نسخه‌ی \lr{JavaScript} \\
\bottomrule
\end{tabular}
\end{table}

ماژول‌های اصلی پایتون عبارت‌اند از: \lr{\texttt{ingest.py}} (تبدیل و بخش‌بندی)، \lr{\texttt{textfix.py}} (ترمیم نویسه‌ها)، \lr{\texttt{index.py}} و \lr{\texttt{qdrant\_store.py}} (نمایه)، \lr{\texttt{retrieve.py}} (بازیابی ترکیبی)، \lr{\texttt{qa.py}} (پاسخ، ترجمه و واژه‌نامه)، \lr{\texttt{extract.py}} و \lr{\texttt{schema.py}} (استخراج ساخت‌یافته)، \lr{\texttt{dispersion.py}} و \lr{\texttt{formulas.py}} (موتور فرمول پاشندگی)، \lr{\texttt{crystal\_names.py}} (جدول نام بلورها) و \lr{\texttt{evaluate.py}} (ارزیابی). ابزار خط فرمان \lr{\texttt{matrag}} همه‌ی این مراحل را با زیرفرمان‌هایی مانند \lr{\texttt{ingest}}، \lr{\texttt{evaluate retrieval}} و \lr{\texttt{reference hitran}} در دسترس می‌گذارد. هر ماژول یک وظیفه‌ی مشخص دارد و از راه ساختارهای داده‌ی ساده (فهرست بخش‌ها با فراداده) با ماژول بعدی ارتباط می‌گیرد، پس هر مرحله را می‌توان جداگانه آزمود یا جایگزین کرد.

فیزیک محاسبه‌گرها عمداً دو بار نوشته شده است. سایت باید بدون سرور و در مرورگر حساب کند (\lr{JavaScript})، و نسخه‌ی پایتون با آزمون‌هایش مرجعی است که عددهای سایت با آن مقایسه می‌شوند. موتور فرمول پاشندگی \lr{JavaScript} با موتور پایتون روی ۶۸۶۷ نقطه مقایسه شد و خروجی‌ها یکسان بودند. این دوبار نوشتن، اگرچه کار بیشتری می‌طلبد، روشی استاندارد برای یافتن خطاهای پیاده‌سازی است: احتمال اینکه دو پیاده‌سازی مستقل خطای یکسانی داشته باشند کم است.

\section{زیرساخت}
جدول \ref{tab:infra} اجزای زیرساخت را نشان می‌دهد. انتخاب‌ها با دو قید انجام شدند: هزینه‌ی صفر یا نزدیک به صفر برای یک پروژه‌ی دانشجویی، و امکان تکرار نتایج.

\begin{table}[htbp]
\centering
\caption{اجزای زیرساخت.}
\label{tab:infra}
\renewcommand{\arraystretch}{1.3}
\small
\begin{tabular}{p{2.8cm}p{4.4cm}p{7.4cm}}
\toprule
جزء & فناوری & نقش \\
\midrule
پردازش سنگین & \lr{Google Colab (GPU T4)} و \lr{GitHub Actions} & تبدیل \lr{PDF}، \lr{OCR} با ۲۰ اجراکننده‌ی موازی، نمایه‌سازی \\
ذخیره‌ی \lr{PDF} & \lr{Google Drive} & آرشیو اصلی و حافظه‌ی نهان تبدیل‌ها؛ تبدیل در قطعه‌های ۴۰ صفحه‌ای \\
پایگاه برداری محلی & \lr{ChromaDB} و \lr{LlamaIndex} & اجرای محلی، آزمون‌ها و ارزیابی \\
پایگاه برداری سایت & \lr{Qdrant Cloud} & ۶۸۱۹۰ بخش با بردار متراکم و \lr{BM25}؛ ادغام \lr{RRF} در خود \lr{Qdrant} \\
سایت و \lr{API} & \lr{Cloudflare Workers} & جست‌وجو، محاسبه‌گرها، احراز هویت، نمایش \lr{PDF} با \lr{pdf.js} \\
مدل‌های زبانی & \lr{Ollama (Qwen3-4B)}، \lr{Gemini}، \lr{Workers AI} & پاسخ و استخراج (محلی)، ترجمه‌ی پرسش فارسی و بردار پرسش (سایت) \\
\bottomrule
\end{tabular}
\end{table}

تبدیل \lr{PDF} در قطعه‌های ۴۰ صفحه‌ای انجام می‌شود تا قطع شدن نشست \lr{Colab} کار انجام‌شده را از بین نبرد؛ خروجی هر قطعه در حافظه‌ی نهان \lr{Google Drive} ذخیره می‌شود و اجرای بعدی از همان‌جا ادامه می‌یابد. متن مقاله‌ها حق نشر دارد. برای همین فایل‌های \lr{PDF} هرگز در مخزن کد قرار نمی‌گیرند و سایت عمومی نیست. سایت سطوح دسترسی جداگانه دارد، از جمله حالتی که فقط ابزارهای محاسبه را بدون نمایش \lr{PDF} در اختیار می‌گذارد.

\section{رابط کاربری و ابزارها}
دو رابط ساخته شد. رابط محلی با \lr{Gradio} این زبانه‌ها را دارد: افزودن مقاله، پرسش و پاسخ (با پاسخ مدل بدون منبع در کنارش و نمایش صفحه‌ی مقاله با بخش استفاده‌شده‌ی رنگی)، بسته‌ی منابع، استخراج خواص، فرمول به جدول ضریب شکست، و رسم طیف جذبی \lr{HITRAN} با \lr{HAPI}. در زبانه‌ی آخر، طیف دوم با مقدارهای استخراج‌شده از مقاله رسم و اختلاف دو طیف نشان داده می‌شود. این مقایسه‌ی دیداری، اثر فیزیکی یک خطای استخراج را مستقیماً نشان می‌دهد.

سایت (\lr{\texttt{matrag.ethanjamescarter1995.workers.dev}}) دوزبانه است و این ابزارها را دارد:
\begin{itemize}
\item \textbf{جست‌وجو:} پرسش فارسی یا انگلیسی؛ نتایج گروه‌بندی‌شده به تفکیک مقاله با نام، سال، صفحه و \lr{DOI}؛ کادر زرد روی صفحه‌ی اصلی \lr{PDF}؛ پاسخ «پیدا نشد». سایت خودش پاسخ نوشتاری تولید نمی‌کند و بسته‌ی منابع را برای مدل‌های زبانی بیرونی آماده می‌کند.
\item \textbf{ضریب شکست:} منحنی $n(\lambda)$، ضریب گروهی و \lr{GVD} برای ۲۳۳۰ ماده، با مقایسه‌ی منابع.
\item \textbf{تطبیق فاز:} فرایندهای \lr{SHG}، \lr{SFG}، \lr{DFG} و \lr{OPO} در ۱۸۶ بلور؛ حل معادله‌ی فرنل برای بلور دومحوری؛ \lr{walk-off}؛ پهنای پذیرش؛ $d_{\sub{eff}}$؛ دمای \lr{NCPM}؛ و پیکربندی‌های \lr{QPM}.
\item \textbf{\lr{SPDC}:} دامنه و شدت طیفی مشترک، خلوص، عدد اشمیت و فرورفتگی \lr{HOM}؛ دمش گاوسی کانونی‌شده (\lr{Bennink 2010} \cite{bennink}).
\item \textbf{مشاور انتخاب بلور:} کاربر شرایط ثابت آزمایشگاه (فرایند، طول موج‌ها، بازه‌ی دمای آون، طول بلور، حداکثر \lr{walk-off} و جز آن) را وارد می‌کند و سامانه همه‌ی ۱۸۶ بلور را می‌سنجد. خروجی فهرست بلورهای مناسب با رتبه و دلیل رد شدن هر بلور دیگر است (۱۶ دلیل تعریف‌شده، مانند جذب در یکی از طول موج‌ها، دمای خارج از بازه‌ی آون یا صفر بودن $d_{\sub{eff}}$).
\end{itemize}

\section{آزمون نرم‌افزار}
تا ۳ آبان ۱۴۰۵، ۲۳۷ آزمون خودکار پایتون و ۸۲ آزمون سمت سایت نوشته شده بود. این آزمون‌ها چهار دسته‌اند: درستی فیزیک (زاویه، دما، دوره‌ی \lr{QPM} و $d_{\sub{eff}}$ در برابر عدد مقاله)، ساختار داده‌ها، کامل بودن ترجمه‌ی انگلیسی رابط، و رفتار \lr{API} در برابر ورودی خراب. آزمون‌های پایتون بدون اینترنت و بدون مدل زبانی اجرا می‌شوند، پس نتیجه‌شان به در دسترس بودن سرویس‌های بیرونی یا رفتار غیرقطعی مدل‌ها وابسته نیست. چند دور اشکال‌یابی روی سایت هم انجام شد و هر اشکال پیداشده یک آزمون تازه به مجموعه اضافه کرد، تا همان اشکال دوباره بازنگردد.

% ======================================== فصل ۵ ======================================================
\chapter{ارزیابی و نتایج}

\section{سنجه‌ها}
کیفیت بازیابی با «پرسش‌های طلایی» سنجیده شد. برای هر پرسش، مقاله‌ی درست (و در مجموعه‌ی آزمایشی، صفحه‌ی درست) از پیش مشخص است. دو سنجه به کار رفت:
\begin{equation}\label{eq:metrics}
\mathrm{R@}k=\frac{1}{|Q|}\sum_{q\in Q}\mathbb{1}\left[\mathrm{rank}_q\le k\right],\qquad
\mathrm{MRR@}10=\frac{1}{|Q|}\sum_{q\in Q}\frac{\mathbb{1}\left[\mathrm{rank}_q\le10\right]}{\mathrm{rank}_q},
\end{equation}
که در آن $\mathrm{rank}_q$ رتبه‌ی نخستین نتیجه‌ی درست برای پرسش $q$ و $\mathbb{1}[\cdot]$ تابع نشانگر است. سنجه‌ی $\mathrm{R@}k$ می‌گوید در چه سهمی از پرسش‌ها پاسخ درست در $k$ نتیجه‌ی اول هست، و $\mathrm{MRR}$ به رتبه‌ی دقیق آن هم حساس است: پاسخ درست در رتبه‌ی اول امتیاز ۱، در رتبه‌ی دوم امتیاز نیم و در رتبه‌ی دهم امتیاز یک‌دهم می‌گیرد. \lr{MRR} برابر ۱ یعنی پاسخ درست همیشه نتیجه‌ی اول است. برای پاسخ‌ها، «دقت» سهم پرسش‌هایی است که پاسخ مقدار درست را دارد. «امتناع درست» هم سهم پرسش‌های بی‌جوابی است که سامانه برایشان «پیدا نشد» گفته است.

برای سنجیدن عدم‌قطعیت آماری، بازه‌ی اطمینان با بازنمونه‌گیری \lr{bootstrap} \cite{efron} محاسبه شد: از مجموعه‌ی $|Q|$ پرسش، $|Q|$ پرسش با جایگذاری برداشته و سنجه دوباره حساب می‌شود؛ پس از ۵۰۰۰ تکرار، صدک‌های ۲/۵ و ۹۷/۵ توزیع حاصل، بازه‌ی اطمینان ۹۵ درصد را می‌دهند. برای مقایسه‌ی دو روش، بازنمونه‌گیری «جفتی» روی همان پرسش‌ها انجام می‌شود، تا تفاوت دشواری پرسش‌ها در هر دو روش یکسان اثر کند.

\section{آزمون اولیه روی مجموعه‌های کوچک}
نخستین ارزیابی روی دو مجموعه‌ی کوچک (سه مقاله در هر حوزه) انجام شد تا درستی خط لوله و ابزارهای ارزیابی سنجیده شود. معیار برخورد در این آزمون سخت‌گیرانه بود: صفحه‌ی درست از مقاله‌ی درست در $k$ بخش اول. جدول \ref{tab:pilot} نتیجه را نشان می‌دهد.

\begin{table}[htbp]
\centering
\caption[بازیابی در مجموعه‌های آزمایشی.]{بازیابی در مجموعه‌های آزمایشی (\lr{\texttt{results/\allowbreak eval/\allowbreak *\_retrieval.csv}}).}
\label{tab:pilot}
\renewcommand{\arraystretch}{1.3}
\begin{tabular}{llcccc}
\toprule
حوزه & روش & \lr{R@1} & \lr{R@3} & \lr{R@5} & \lr{MRR} \\
\midrule
\multirow{3}{*}{\lr{HITRAN} (۱۵ پرسش)} & برداری & ۰/۶۰ & ۰/۷۳ & ۰/۸۷ & ۰/۷۲ \\
 & \lr{BM25} & ۰/۶۰ & ۰/۸۰ & ۰/۸۷ & ۰/۷۲ \\
 & ترکیبی & ۰/۷۳ & ۰/۸۰ & ۰/۸۷ & ۰/۷۹ \\
\midrule
\multirow{3}{*}{اپتیک (۱۳ پرسش)} & برداری & ۰/۶۹ & ۱/۰۰ & ۱/۰۰ & ۰/۸۵ \\
 & \lr{BM25} & ۰/۹۲ & ۱/۰۰ & ۱/۰۰ & ۰/۹۶ \\
 & ترکیبی & ۰/۶۹ & ۱/۰۰ & ۱/۰۰ & ۰/۸۳ \\
\bottomrule
\end{tabular}
\end{table}

در \lr{HITRAN}، بازیابی ترکیبی بهترین نتیجه را داد. در مجموعه‌ی اپتیک، \lr{BM25} به‌تنهایی بهتر بود. علت احتمالی این است که پرسش‌های اپتیک نام ماده و نماد دقیق دارند و در مجموعه‌ای با سه مقاله، تطابق کلیدواژه تقریباً همیشه کافی است. در این حالت، افزودن فهرست برداری فقط نامزدهای نادرستی را به بالای فهرست ادغام‌شده می‌آورد.

\subsection{\lr{RAG} در برابر مدل بدون منبع}
جدول \ref{tab:hitran-answers} پاسخ سامانه را با پاسخ همان مدل بدون منبع مقایسه می‌کند.

\begin{table}[htbp]
\centering
\caption[پاسخ به پرسش‌های آزمون \lr{HITRAN}.]{پاسخ به ۱۵ پرسش جواب‌دار و ۱ پرسش بی‌جواب \lr{HITRAN} (\lr{\texttt{results/\allowbreak eval/\allowbreak hitran\_answers.csv}}).}
\label{tab:hitran-answers}
\renewcommand{\arraystretch}{1.3}
\resizebox{\textwidth}{!}{%
\begin{tabular}{lccc}
\toprule
سامانه & دقت (پرسش جواب‌دار) & «پیدا نشد» برای پرسش جواب‌دار & امتناع درست (پرسش بی‌جواب) \\
\midrule
\lr{RAG + Qwen3-4B} & ۰/۹۳ (۱۴ از ۱۵) & ۰/۰۰ & ۱/۰ \\
\lr{Qwen3-4B} بدون منبع & ۰/۰۰ & ۱/۰۰ & ۱/۰ \\
\bottomrule
\end{tabular}}
\end{table}

مدل کوچک بدون منبع در هیچ پرسشی عدد درست را نیاورد و همیشه گفت پاسخ را نمی‌داند. این رفتار از تولید عدد نادرست بهتر است، اما کاربرد عملی ندارد. با \lr{RAG}، همان مدل در ۱۴ پرسش از ۱۵ پرسش مقدار درست را با منبع آورد. تنها خطا (پرسش \lr{h07}) از نوع «دقت عددی» بود. منبع درست بازیابی شده بود، اما مدل به‌جای مقدار دقیق $5007.787078\ \mathrm{cm^{-1}}$ عدد تقریبی «حدود ۵۰۰۰ $\mathrm{cm^{-1}}$» را از مقدمه‌ی مقاله آورد.

\subsection{نمونه‌ی پاسخ‌ها}
جدول \ref{tab:samples} چند نمونه از پاسخ‌های واقعی ثبت‌شده در این آزمون را نشان می‌دهد (\lr{\texttt{results/\allowbreak eval/\allowbreak hitran\_answers\_details.csv}}). پرسش‌ها به انگلیسی پرسیده شده‌اند و در اینجا ترجمه شده‌اند. پاسخ‌ها بی‌تغییر آمده‌اند.

\begin{table}[htbp]
\centering
\caption{نمونه‌ی پاسخ‌های ثبت‌شده در آزمون \lr{HITRAN}.}
\label{tab:samples}
\renewcommand{\arraystretch}{1.3}
\small
\begin{tabular}{p{6.0cm}p{5.4cm}p{3.2cm}}
\toprule
پرسش & پاسخ \lr{RAG + Qwen3-4B} & مدل بدون منبع \\
\midrule
بر پایه‌ی محاسبات نظری، پهنای نیم‌بیشینه‌ی پهن‌شدگی خطوط $\mathrm{CO_2}$ با بخار آب در مقایسه با پهن‌شدگی با هوا تا چه اندازه می‌تواند بزرگ شود؟ & \lr{\ldots can reach twice the values of the air-broadened half-widths [1].} & \lr{Not found in the provided sources.} \\
چه ضریب مقیاسی روی پارامترهای پهن‌شدگی با هوای $\mathrm{CH_4}$ اعمال شد تا ضرایب پهن‌شدگی با بخار آب به دست آید؟ & \lr{\ldots was 1.36 [1].} & \lr{Not found in the provided sources.} \\
\lr{Sung} و همکاران پهن‌شدگی با بخار آب را برای چند خط $^{12}\mathrm{CO_2}$ گزارش کردند؟ & \lr{182 lines \ldots in the $\nu_3$ and $\nu_2+\nu_3-\nu_2$ bands [1].} & \lr{Not found in the provided sources.} \\
بازه‌ی ضرایب پهن‌شدگی $\mathrm{H_2}$ که \lr{Padmanabhan} و همکاران برای خطوط باند \lr{30012$\leftarrow$00001} اندازه گرفتند چیست؟ & \lr{\ldots from 0.095 cm$^{-1}$/atm for the P(32) line to 0.121 cm$^{-1}$/atm for the P(20) line [1].} & \lr{Not found in the provided sources.} \\
ضریب کالیبراسیون مقیاس عدد موج طیف‌های تبدیل فوریه‌ی $\mathrm{CO_2}$ و $\mathrm{H_2}$ چقدر بود؟ & \lr{\ldots was 1.000000467 [1].} & \lr{Not found in the provided sources.} \\
ضریب پهن‌شدگی خودی خط $\mathrm{NH_3}$ در ۱/۲ میکرومتر در دمای ۱۵۰ کلوین چقدر است؟ (پرسش بی‌جواب) & \lr{Not found in the provided sources.} & \lr{Not found in the provided sources.} \\
\bottomrule
\end{tabular}
\end{table}

دو نکته در این جدول دیده می‌شود. پاسخ \lr{RAG} عدد را با همان دقت منبع می‌آورد، حتی وقتی عدد ۱۰ رقم معنادار دارد (\lr{1.000000467})، و شرایط اندازه‌گیری مانند نام خط و باند را هم حفظ می‌کند. پرسش آخر عمداً بی‌جواب طراحی شده بود: مقاله‌ای درباره‌ی $\mathrm{NH_3}$ در مجموعه نبود و سامانه به‌جای ساختن عدد، «پیدا نشد» گفت.

\subsection{اثر پرسش فارسی}
همان ۱۵ پرسش به فارسی نوشته و پس از ترجمه‌ی خودکار جست‌وجو شدند (جدول \ref{tab:fa}).

\begin{table}[htbp]
\centering
\caption[اثر پرسش فارسی بر بازیابی.]{همان ۱۵ پرسش \lr{HITRAN} به فارسی، ترجمه با \lr{Qwen3-4B} (\lr{\texttt{results/\allowbreak eval/\allowbreak hitran\_retrieval\_fa.csv}}).}
\label{tab:fa}
\renewcommand{\arraystretch}{1.3}
\begin{tabular}{lcc}
\toprule
روش & \lr{MRR} پرسش انگلیسی & \lr{MRR} پرسش فارسی (پس از ترجمه) \\
\midrule
برداری & ۰/۷۲ & ۰/۷۲ \\
\lr{BM25} & ۰/۷۲ & ۰/۶۳ \\
ترکیبی & ۰/۷۹ & ۰/۷۲ \\
\bottomrule
\end{tabular}
\end{table}

جست‌وجوی برداری از ترجمه اثر نپذیرفت، ولی \lr{BM25} که به واژه‌ی دقیق وابسته است افت کرد. این تفاوت با سازوکار دو روش سازگار است: ترجمه معنای پرسش را حفظ می‌کند ولی واژه‌ی دقیقی را که نویسنده‌ی مقاله به کار برده لزوماً تولید نمی‌کند. خطای شاخص ترجمه ابهام واژه‌ی «چند» بود: «چند خط» در دو پرسش به «\lr{several lines}» ترجمه شد، در حالی که منظور «\lr{how many}» بود.

این عددها روی مجموعه‌هایی با ۱۳ تا ۱۶ پرسش و سه مقاله به دست آمده‌اند. آن‌ها فقط نشان می‌دهند خط لوله و ابزار ارزیابی درست کار می‌کنند و برای نتیجه‌گیری کلی کافی نیستند. با ۱۵ پرسش، تغییر پاسخ تنها یک پرسش سنجه‌ی $\mathrm{R@}k$ را حدود ۷ درصد جابه‌جا می‌کند.

\section{ارزیابی جست‌وجو روی مجموعه‌ی بزرگ}
برای مجموعه‌ی ۲۸۴۵ سندی، بانک پرسش‌ها از سه منبع ساخته شد: پرسش‌های ساخته‌شده از داده‌های \lr{refractiveindex.info} (گروه \lr{ri}: «فرمول ضریب شکست ماده‌ی \lr{X} از کدام مقاله است؟»)، پرسش‌هایی که مدل زبانی از متن مقاله‌ها نوشت (گروه \lr{gen})، و پرسش‌هایی از داده‌های آزموده‌ی پایگاه (گروه‌های \lr{dij}، \lr{thermo} و \lr{paper}). ارزیابی روی سایت زنده اجرا شد. معیار برخورد، حضور مقاله‌ی طلایی در ۱۰ نتیجه‌ی اول است. جدول \ref{tab:big} نتیجه را نشان می‌دهد.

\begin{table}[htbp]
\centering
\caption[ارزیابی ۲۰۰ پرسش طلایی روی سایت زنده.]{ارزیابی ۲۰۰ پرسش طلایی روی سایت زنده، ۴ آبان ۱۴۰۵ (\lr{\texttt{scripts/\allowbreak quality/\allowbreak reports/\allowbreak search-2026-10-26.md}}). بازه‌ی اطمینان با بازنمونه‌گیری \lr{bootstrap} (۵۰۰۰ بار) روی پرسش‌ها محاسبه شد.}
\label{tab:big}
\renewcommand{\arraystretch}{1.3}
\resizebox{\textwidth}{!}{%
\begin{tabular}{lccccc}
\toprule
روش & \lr{R@1} & \lr{R@3} & \lr{R@10} & \lr{MRR@10} & بازه‌ی اطمینان ۹۵ درصد \\
\midrule
جست‌وجوی ترکیبی (ترتیب سایت) & ۰/۴۳۰ & ۰/۵۹۵ & ۰/۶۹۵ & ۰/۵۱۸ & ۰/۴۶ تا ۰/۵۸ \\
رتبه‌بندی دوم \lr{SubQ} روی ۳۰ نامزد & ۰/۵۰۰ & ۰/۶۵۵ & ۰/۷۴۵ & ۰/۵۹۰ & ۰/۵۳ تا ۰/۶۵ \\
\bottomrule
\end{tabular}}
\end{table}

بهبود حاصل از رتبه‌بندی دوم برابر $+0.071$ است. بازه‌ی اطمینان ۹۵ درصد جفتی این تفاوت (بازنمونه‌گیری روی همان پرسش‌ها) از $+0.027$ تا $+0.115$ است و صفر را دربر نمی‌گیرد. پس این بهبود از نظر آماری معنادار است. توجه کنید که دو بازه‌ی جدول \ref{tab:big} با هم هم‌پوشانی دارند، اما این به معنای نبود تفاوت معنادار نیست: بازه‌ی جفتی دقیق‌تر است، چون دشواری پرسش‌ها در هر دو روش یکسان است و از تفاوت حذف می‌شود.

جدول \ref{tab:groups} و شکل \ref{fig:mrr} همین نتیجه را به تفکیک گروه پرسش نشان می‌دهند.

\begin{table}[htbp]
\centering
\caption{\lr{MRR@10} به تفکیک گروه پرسش (همان ارزیابی جدول \ref{tab:big}).}
\label{tab:groups}
\renewcommand{\arraystretch}{1.3}
\begin{tabular}{lccc}
\toprule
گروه & تعداد & ترکیبی & با \lr{SubQ} \\
\midrule
\lr{gen}: پرسش نوشته‌شده از متن مقاله & ۷۶ & ۰/۷۵۴ & ۰/۷۸۱ \\
\lr{ri}: مقاله‌ی منبع فرمول ضریب شکست & ۱۱۰ & ۰/۳۷۴ & ۰/۴۸۱ \\
\lr{dij}: ضرایب غیرخطی & ۷ & ۰/۳۸۶ & ۰/۴۸۲ \\
\lr{thermo}: ضریب ترمواپتیک & ۴ & ۰/۳۷۵ & ۰/۲۷۱ \\
\lr{paper}: مقاله‌ی فرمول & ۳ & ۰/۳۳۳ & ۰/۴۱۷ \\
\midrule
کل & ۲۰۰ & ۰/۵۱۸ & ۰/۵۹۰ \\
\bottomrule
\end{tabular}
\end{table}

\begin{figure}[htbp]
\centering
\figbox{fig_mrr_groups.png}{0.85\textwidth}
\caption{\lr{MRR@10} به تفکیک گروه پرسش در ارزیابی ۲۰۰ پرسشی؛ ستون کم‌رنگ جست‌وجوی ترکیبی و ستون پررنگ با رتبه‌بندی دوم \lr{SubQ}. عدد داخل پرانتز تعداد پرسش‌های هر گروه است.}
\label{fig:mrr}
\end{figure}

پرسش‌هایی که از متن خود مقاله نوشته شده‌اند (\lr{gen}) نتیجه‌ی خوبی دارند. این انتظار می‌رود، چون واژه‌های این پرسش‌ها از همان متن آمده‌اند. ضعیف‌ترین گروه بزرگ، \lr{ri} است. در این گروه فقط مقاله‌ای درست حساب می‌شود که \lr{refractiveindex.info} به آن ارجاع داده است، در حالی که اغلب چند مقاله‌ی دیگر مجموعه هم ضریب شکست همان ماده را اندازه گرفته‌اند. برای نمونه در پرسش $\mathrm{YVO_4}$، مقاله‌های \lr{Zelmon 2010} و \lr{Zeng 1990} که هر دو اندازه‌گیری ضریب شکست $\mathrm{YVO_4}$ هستند بالاتر از مقاله‌ی طلایی \lr{Lomheim 1978} آمدند. پس این سنجه عملکرد واقعی را کمتر از واقع نشان می‌دهد. گروه‌های \lr{dij}، \lr{thermo} و \lr{paper} فقط ۳ تا ۷ پرسش دارند و عددشان برای نتیجه‌گیری کافی نیست؛ برای نمونه، افت گروه \lr{thermo} با رتبه‌بندی دوم از تغییر رتبه‌ی یکی دو پرسش ناشی می‌شود.

\subsection{روند در طول پروژه}
در گزارش اول به استاد راهنما (۱۳ مهر ۱۴۰۵) مجموعه ۲۶۲۳ مقاله و ۵۶۹۱۳ بخش داشت و \lr{MRR} کل روی بانک ۷۱۷ پرسشی ۰/۴۵۵ بود. در یک ارزیابی دروازه‌ی کیفیت با ۸۰ پرسش روی سایت زنده (۲ آبان ۱۴۰۵)، \lr{MRR@10} جست‌وجوی ترکیبی ۰/۵۳۵ و با رتبه‌بندی دوم ۰/۶۱۳ بود. این عددها روی مجموعه‌های پرسش متفاوتی به دست آمده‌اند و مستقیماً با هم قابل‌مقایسه نیستند.

\section{آزمایش‌های مقایسه‌ای}\label{sec:ablation}
اثر هر تغییر، پیش از پذیرفته شدن در سامانه، به‌صورت کمّی اندازه‌گیری شد. تغییرهایی که نتیجه را بهتر نکردند، هرچند از نظر نظری موجه بودند، کنار گذاشته شدند. جدول \ref{tab:ablation} مهم‌ترین این آزمایش‌ها را نشان می‌دهد.

\begin{table}[htbp]
\centering
\caption{تغییرات مهم و اثر اندازه‌گیری‌شده‌ی آن‌ها.}
\label{tab:ablation}
\renewcommand{\arraystretch}{1.3}
\small
\begin{tabular}{p{7.0cm}p{1.6cm}p{1.6cm}p{3.0cm}}
\toprule
تغییر & قبل & بعد & تصمیم \\
\midrule
رتبه‌بند متقاطع \lr{bge-reranker-base} روی ۳۰ نامزد (۲۰۰ پرسش، ۱۸ مهر) & ۰/۴۳۹ & ۰/۴۰۱ & کنار گذاشته شد \\
حذف بخش‌های کم‌ارزش با نقشه‌ی داده & ۰/۵۱۱ & ۰/۴۸۵ & کنار گذاشته شد \\
جست‌وجوی جداگانه روی نام ماده: تعداد «پیدا نشد»های اشتباه & ۵۲ & ۱۳ & نگه داشته شد \\
پیش‌واکشی متمرکز بر بلور (\lr{MRR}) & ۰/۵۱۸ & ۰/۵۳۸ & نگه داشته شد \\
رتبه‌بندی دوم \lr{SubQ} (ارزیابی ۸۰ پرسشی) & ۰/۵۳۲ & ۰/۶۱۴ & پیش‌فرض شد \\
بهینه‌سازی سرعت: زمان پاسخ جست‌وجو (ثانیه) & ۲/۱ تا ۲/۹ & ۱/۱ تا ۱/۴ & کیفیت بدون تغییر ماند \\
\bottomrule
\end{tabular}
\end{table}

دو نکته در تفسیر این جدول مهم است. اول، رتبه‌بند \lr{bge} در گروه \lr{ri} افت شدید داشت (۰/۱۹۷ به ۰/۱۴۴) و در گروه \lr{gen} تقریباً بی‌اثر بود (۰/۸۲۱ به ۰/۸۱۹). یک توضیح محتمل این است که این رتبه‌بند روی متن عمومی آموزش دیده و پرسش‌های گروه \lr{ri}، که فقط نام ماده و «ضریب شکست» دارند، برایش اطلاعات کافی برای تمایز میان مقاله‌های هم‌موضوع فراهم نمی‌کنند. دوم، با توجه به بازه‌ی اطمینان حدود $\pm0.06$ در جدول \ref{tab:big}، تفاوت‌های حدود ۰/۰۲ (مانند پیش‌واکشی بلور) در حد نوسان آماری‌اند و نمی‌توان با اطمینان گفت بهبود واقعی‌اند. افزون بر این، این تنظیم‌ها روی همان پرسش‌هایی انجام شد که گزارش هم روی آن‌هاست، و این می‌تواند نتیجه را خوش‌بینانه کند.

\section{درستی‌آزمایی فیزیکی}
خروجی محاسبه‌گرها با عددهایی مقایسه شد که خود مقاله‌ها گزارش کرده‌اند (جدول \ref{tab:physics}). همه‌ی این موارد در کد به‌صورت آزمون خودکار وجود دارند.

\begin{table}[htbp]
\centering
\caption{خروجی سامانه در برابر عدد گزارش‌شده در مقاله‌ی اصلی.}
\label{tab:physics}
\renewcommand{\arraystretch}{1.3}
\small
\begin{tabular}{p{7.4cm}p{3.6cm}p{3.6cm}}
\toprule
آزمون & سامانه & مقاله \\
\midrule
\lr{BBO}، زاویه‌ی \lr{SHG} نوع اول در ۱/۰۶۴ میکرومتر & ۲۲/۸۳ درجه & ۲۲/۸ درجه \\
\lr{GdCOB}، زاویه در صفحه‌های \lr{xz} و \lr{xy} (\lr{Aka 1997}) & ۱۹/۶۸ و ۴۵/۹۹ درجه & ۱۹/۶۸ و ۴۵/۹۹ درجه \\
$\mathrm{CdSiP_2}$، \lr{SHG} در ۴/۷۸۴۶ میکرومتر (\lr{Kato 2011}) & ۴۳/۱۴ درجه & ۴۳/۱۳ درجه \\
$d_{\sub{eff}}$ بلورهای \lr{BBO} و \lr{LBO} (پیکومتر بر ولت) & ۲/۰۲ و ۰/۸۳ & ۱/۹۴ و ۰/۸۲ \\
$d_{\sub{eff}}$ بلورهای $\mathrm{LiInS_2}$ و $\mathrm{LiInSe_2}$ (\lr{Petrov 2010}) & ۶/۵۱ و ۹/۳۴ & ۶/۵۴ و ۹/۳۵ \\
دوره‌ی \lr{QPM} بلورهای \lr{PPLN} و \lr{PPKTP} برای \lr{SHG} ۱/۰۶۴ میکرومتر (میکرومتر) & ۶/۹ و ۹/۰ & حدود ۶/۸ تا ۶/۹ و ۹/۰ \\
\lr{PPKTP}، \lr{SPDC} با دمش ۴۰۵ نانومتر (\lr{Fedrizzi 2007} \cite{fedrizzi}) & ۱۰/۰ میکرومتر در ۴۹/۲ درجه‌ی سانتی‌گراد & ۱۰/۰ میکرومتر، حدود ۴۹ درجه \\
\lr{KTP}، دمای \lr{NCPM} سه فرایند (\lr{Kato 2002}) (درجه‌ی سانتی‌گراد) & ۶۳/۷۷ و ۷۷/۵۶ و ۵۹/۶۴ & ۶۳/۸ و ۷۷/۶ و ۵۹/۸ \\
\lr{LBO}، دمای \lr{NCPM} برای \lr{SHG} ۱/۰۶۴ میکرومتر (درجه‌ی سانتی‌گراد) & حدود ۱۴۹ & ۱۴۸ تا ۱۴۹/۵ (اندازه‌گیری) \\
$\mathrm{CaF_2}$ (\lr{Leviton} و همکاران)، ضریب شکست در ۰/۴۰ میکرومتر و ۲۹۵ کلوین & ۱/۴۴۲۲۱ & بازتولید تا رقم پنجم \\
\bottomrule
\end{tabular}
\end{table}

آزمون منظم داده‌ها چند ناسازگاری را در خود ادبیات یا در نقل آن آشکار کرد: یک عدد تکراری در جدول ضرایب \lr{CBO}، اختلاف ۰/۰۳ در $n_x$ فرمول \lr{YCOB} نسبت به جدول اندازه‌گیری همان مقاله، جابه‌جایی $d_{31}$ و $d_{32}$ در \lr{KTP} و \lr{LBO} میان قراردادهای مختلف محور، و نگاشت محورهای بلورشناسی به دی‌الکتریک در $\mathrm{LiInS_2}$، $\mathrm{LiInSe_2}$، \lr{BiBO} و \lr{MNA}. این یافته‌ها نشان می‌دهند که آزمون نظام‌مند داده‌ها در برابر عدد مقاله، افزون بر کنترل کیفیت، خود ارزش علمی مستقلی دارد.

مشاور انتخاب بلور با ۱۵ سناریوی نمونه آزموده شد و در همه‌ی آن‌ها پیشنهادهای اول با انتخاب رایج متخصصان هم‌خوان بود: \lr{BBO} و \lr{CLBO} برای تولید فرابنفش، \lr{PPLN} و \lr{PPKTP} برای تبدیل در ناحیه‌ی مرئی، و \lr{GaSe}، $\mathrm{ZnGeP_2}$، $\mathrm{CdSiP_2}$ و \lr{OP-GaAs} برای فروسرخ میانی. این سناریوها را خود نگارنده ساخته است و آزمون با سناریوهای واقعی آزمایشگاه‌ها هنوز انجام نشده است.

\subsection{مثال حل‌شده: تولید هماهنگ دوم نوع اول در \lr{BBO}}
برای نشان دادن زنجیره‌ی محاسبه از داده تا نتیجه، تولید هماهنگ دوم لیزر \lr{Nd:YAG} ($\lambda_\omega=1.064\ \mu\mathrm{m}$) در بلور تک‌محوری منفی \lr{BBO} به‌طور کامل حل می‌شود. فرمول سلمایر به‌کاررفته همان فرمولی است که در بسته‌ی مرجع پایتون پروژه برای \lr{BBO} ثبت شده و از کار \lr{Eimerl} و همکاران گرفته شده است \cite{eimerl}:
\begin{equation}\label{eq:bbo}
\begin{aligned}
n_o^{2}&=2.7359+\frac{0.01878}{\lambda^{2}-0.01822}-0.01354\,\lambda^{2},\\
n_e^{2}&=2.3753+\frac{0.01224}{\lambda^{2}-0.01667}-0.01516\,\lambda^{2}.
\end{aligned}
\end{equation}
همه‌ی عددهای این بخش و بخش بعد با برنامه‌ی \lr{\texttt{scripts/\allowbreak calc\_figures.py}} (پیوست ج) دوباره محاسبه شدند و با عددهای ثبت‌شده در بسته‌ی مرجع پروژه یکسان‌اند. شکل \ref{fig:bbo-n} منحنی دو ضریب شکست را در ناحیه‌ی شفافیت نشان می‌دهد.

\begin{figure}[htbp]
\centering
\figbox{fig_bbo_dispersion.png}{0.78\textwidth}
\caption{ضریب شکست عادی ($n_o$) و غیرعادی ($n_e$) بلور \lr{BBO} از رابطه‌ی (\ref{eq:bbo}). خط‌چین‌های عمودی طول موج‌های ۰/۵۳۲ و ۱/۰۶۴ میکرومتر را نشان می‌دهند. دوشکستی منفی ($n_e<n_o$) شرط لازم تطبیق فاز نوع اول است.}
\label{fig:bbo-n}
\end{figure}

\textbf{گام ۱، ضرایب شکست.} با $\lambda$ بر حسب میکرومتر: $n_o(1.064)=1.65452$، $n_o(0.532)=1.67421$ و $n_e(0.532)=1.55466$. به دلیل پاشندگی عادی، $n_o(0.532)>n_o(1.064)$ است و بدون استفاده از دوشکستی تطبیق فاز ممکن نیست.

\textbf{گام ۲، زاویه‌ی تطبیق فاز.} در برهم‌کنش نوع اول (\lr{ooe}) باید $n_e(2\omega,\theta)=n_o(\omega)$ باشد. با ضریب شکست غیرعادی وابسته به زاویه
\begin{equation}\label{eq:theta}
\frac{1}{n_e^{2}(\theta)}=\frac{\cos^{2}\theta}{n_o^{2}}+\frac{\sin^{2}\theta}{n_e^{2}}
\;\Longrightarrow\;
\sin^{2}\theta_{pm}=\frac{n_o(\omega)^{-2}-n_o(2\omega)^{-2}}{n_e(2\omega)^{-2}-n_o(2\omega)^{-2}},
\end{equation}
جای‌گذاری عددها $\theta_{pm}=22.78^\circ$ را می‌دهد. سایت با فرمول پاشندگی خودش $22.83^\circ$ و ادبیات $22.8^\circ$ را گزارش می‌کنند. اختلاف حدود $0.05^\circ$ از تفاوت فرمول‌های سلمایر منابع مختلف می‌آید و نشان می‌دهد که انتخاب منبع داده، حتی برای بلور جاافتاده‌ای مانند \lr{BBO}، در رقم دوم اعشار زاویه اثر دارد.

\textbf{گام ۳، طول همدوسی بدون تطبیق فاز.} اگر نور در امتداد محور نوری ($\theta=0$) منتشر شود، $\Delta k=(2\pi/\lambda_{2\omega})\,[n_o(2\omega)-n_o(\omega)]=0.233\ \mu\mathrm{m^{-1}}$ و از رابطه‌ی (\ref{eq:lc}) $L_c\approx13.5\ \mu\mathrm{m}$. یعنی بدون تطبیق فاز، فقط چند ده میکرومتر از بلور در تولید هماهنگ دوم مؤثر است، در حالی که با تطبیق فاز طول مؤثر به چند میلی‌متر می‌رسد.

\textbf{گام ۴، جدایی مسیر (\lr{walk-off}).} در بلور ناهمسانگرد، جهت انرژی موج غیرعادی با بردار موج زاویه‌ی $\rho$ می‌سازد:
\begin{equation}\label{eq:walkoff}
\tan(\theta+\rho)=\left(\frac{n_o}{n_e}\right)^{2}\tan\theta ,
\end{equation}
که برای موج هماهنگ دوم $\rho=3.19^\circ$ (۵۵/۶ میلی‌رادیان) را می‌دهد. برای باریکه‌ای با قطر ۱۰۰ میکرومتر، دو باریکه پس از حدود ۱/۸ میلی‌متر از هم جدا می‌شوند. همین کمیت طول مفید بلور را برای باریکه‌های کانونی‌شده محدود می‌کند.

\textbf{گام ۵، پهنای پذیرش زاویه‌ای.} مشتق عددی $\Delta k$ نسبت به $\theta$ در زاویه‌ی تطبیق فاز $|\partial\Delta k/\partial\theta|=1.09\ \mu\mathrm{m^{-1}\,rad^{-1}}$ است و از رابطه‌ی (\ref{eq:acc}) $\Delta\theta\cdot L\approx0.51\ \mathrm{mrad\cdot cm}$ به دست می‌آید. پس برای بلوری به طول ۱ سانتی‌متر، زاویه‌ی برش و تنظیم باید بهتر از حدود نیم میلی‌رادیان باشد. شکل \ref{fig:bbo-acc} منحنی کامل $\sinc^2(\Delta kL/2)$ را بر حسب انحراف زاویه برای همین بلور نشان می‌دهد.

\begin{figure}[htbp]
\centering
\figbox{fig_bbo_acceptance.png}{0.78\textwidth}
\caption{توان بهنجار هماهنگ دوم بر حسب انحراف از زاویه‌ی تطبیق فاز در بلور \lr{BBO} به طول ۱ سانتی‌متر (\lr{SHG} نوع اول در ۱/۰۶۴ میکرومتر)، محاسبه‌شده از رابطه‌های (\ref{eq:shg}) و (\ref{eq:bbo}). خط‌چین افقی نصف بیشینه را نشان می‌دهد؛ پهنای کامل منحنی در این تراز حدود ۰/۵ میلی‌رادیان است.}
\label{fig:bbo-acc}
\end{figure}

\textbf{گام ۶، ضریب مؤثر.} برای گروه نقطه‌ای $3m$ در برهم‌کنش \lr{ooe}، $d_{\sub{eff}}=d_{31}\sin\theta-d_{22}\cos\theta\sin3\varphi$. با مقادیر پایگاه ($|d_{22}|=2.2$ و $d_{31}=0.1$ پیکومتر بر ولت) و $\varphi=90^\circ$، جمله‌ی غالب $d_{22}\cos\theta=2.03\ \mathrm{pm/V}$ است و سهم $d_{31}\sin\theta$ فقط $0.04\ \mathrm{pm/V}$. سایت $2.02\ \mathrm{pm/V}$ را محاسبه می‌کند و مقدار اندازه‌گیری‌شده‌ی \lr{Eckardt} و همکاران $1.94\ \mathrm{pm/V}$ است. جدول \ref{tab:bbo} نتیجه‌ها را خلاصه می‌کند.

\begin{table}[htbp]
\centering
\caption{خلاصه‌ی مثال حل‌شده‌ی \lr{SHG} نوع اول در \lr{BBO} در ۱/۰۶۴ میکرومتر.}
\label{tab:bbo}
\renewcommand{\arraystretch}{1.3}
\begin{tabular}{lcc}
\toprule
کمیت & مقدار محاسبه‌شده & مقایسه \\
\midrule
زاویه‌ی تطبیق فاز $\theta_{pm}$ & ۲۲/۷۸ درجه & سایت ۲۲/۸۳، ادبیات ۲۲/۸ درجه \\
طول همدوسی در $\theta=0$ & ۱۳/۵ میکرومتر & \lr{-} \\
زاویه‌ی \lr{walk-off} & ۳/۱۹ درجه (۵۵/۶ میلی‌رادیان) & \lr{-} \\
پهنای پذیرش زاویه‌ای $\Delta\theta\cdot L$ & ۰/۵۱ میلی‌رادیان سانتی‌متر & \lr{-} \\
$d_{\sub{eff}}$ & ۲/۰۳ پیکومتر بر ولت & سایت ۲/۰۲، اندازه‌گیری ۱/۹۴ \\
\bottomrule
\end{tabular}
\end{table}

\subsection{مثال حل‌شده: تطبیق فاز شبه‌فازی در \lr{PPLN}}
در لیتیم نیوبات قطبیده‌ی دوره‌ای، هر سه موج با قطبش غیرعادی منتشر می‌شوند تا از بزرگ‌ترین ضریب، $d_{33}$، استفاده شود. در این هندسه تطبیق فاز دوشکستی ممکن نیست و ناهمخوانی فاز با وارون کردن دوره‌ای قطبش جبران می‌شود. برای \lr{SHG} دوره‌ی لازم برابر است با
\begin{equation}\label{eq:ppln}
\Lambda=\frac{2\pi}{\Delta k}=\frac{\lambda_\omega}{2\,[n_e(2\omega)-n_e(\omega)]}.
\end{equation}
با فرمول وابسته به دمای \lr{Gayer} و همکاران برای لیتیم نیوبات آلاییده با \lr{MgO} \cite{gayer}، در دمای ۲۵ درجه‌ی سانتی‌گراد $n_e(1.064)=2.14829$ و $n_e(0.532)=2.22458$ است و $\Lambda=6.97\ \mu\mathrm{m}$ به دست می‌آید. در دمای ۱۰۰ درجه، دوره به $6.84\ \mu\mathrm{m}$ می‌رسد. شکل \ref{fig:ppln} این وابستگی را در بازه‌ی ۲۰ تا ۲۰۰ درجه نشان می‌دهد. این وابستگی دمایی همان چیزی است که در عمل برای تنظیم دقیق تطبیق فاز با دمای آون به کار می‌رود: بلور با دوره‌ی ثابت ساخته می‌شود و دمای آون طوری تنظیم می‌شود که دوره‌ی لازم با دوره‌ی ساخته‌شده یکی شود. مقدار گزارش‌شده‌ی سایت $6.9\ \mu\mathrm{m}$ و مقدار رایج ادبیات $6.8$ تا $6.9\ \mu\mathrm{m}$ است.

\begin{figure}[htbp]
\centering
\figbox{fig_ppln_period.png}{0.78\textwidth}
\caption{دوره‌ی تطبیق فاز شبه‌فازی مرتبه‌ی اول برای تولید هماهنگ دوم ۱/۰۶۴ میکرومتر در لیتیم نیوبات آلاییده با \lr{MgO}، بر حسب دما، از رابطه‌ی (\ref{eq:ppln}) و فرمول \lr{Gayer} و همکاران.}
\label{fig:ppln}
\end{figure}

ضریب مؤثر از رابطه‌ی (\ref{eq:qpm}) با مقدار $|d_{33}|=25.0\ \mathrm{pm/V}$ که پایگاه از کار \lr{Shoji} و همکاران ثبت کرده است، برابر $d_{\sub{eff}}=(2/\pi)|d_{33}|\approx15.9\ \mathrm{pm/V}$ است، یعنی نزدیک به هشت برابر \lr{BBO}. چون بازده تبدیل با $d_{\sub{eff}}^2$ متناسب است (رابطه‌ی (\ref{eq:shg}))، این به معنای بازده‌ای حدود ۶۰ برابر در طول بلور یکسان است، اگر از تفاوت ضرایب شکست در مخرج رابطه چشم‌پوشی شود. این مقایسه دلیل اصلی کاربرد گسترده‌ی \lr{PPLN} در تبدیل طول موج در ناحیه‌ی مرئی و فروسرخ نزدیک است.

\section{استخراج خودکار فرمول: یک نتیجه‌ی منفی}\label{sec:negative}
استخراج خودکار ضرایب فرمول پاشندگی با مدل زبانی روی ۲۴ مقاله آزموده شد و ضریب شکست حاصل با مدخل همان مقاله در \lr{refractiveindex.info} مقایسه شد. در هیچ موردی ضریب شکست محاسبه‌شده تا دقت ۰/۰۱ با مرجع یکی نبود. فرمول پاشندگی به همه‌ی ضرایب، توان‌ها، علامت‌ها و واحد طول موج حساس است و یک خطای کوچک در هر کدام کل منحنی را جابه‌جا می‌کند. برای نمونه، اگر $C_i$ در رابطه‌ی (\ref{eq:sellmeier}) به‌جای مربع طول موج تشدید، خود طول موج تشدید خوانده شود، یا طول موج به‌جای میکرومتر بر حسب نانومتر وارد شود، منحنی حاصل از نظر شکل هم با منحنی درست تفاوت دارد. این نتیجه شیوه‌ی کار بعدی پروژه را تعیین کرد: استخراج نیمه‌دستی و سپس آزمون با عدد مقاله (بخش ۳-۷).

در یک ممیزی خودکار بعدی، مقدار هر داده‌ی محاسبه‌گر کنار بخش منبعش گذاشته و با مدل زبانی سنجیده شد. ده موردی که کمترین امتیاز را گرفتند دستی بازبینی شدند و هیچ عدد اشتباهی پیدا نشد. علت امتیاز پایین، جدول‌های خراب یا تصویری و نگاشت عمدی محورها بود. این نتیجه نشان می‌دهد که مدل زبانی برای «بازبینی» یک عدد موجود در کنار منبعش ابزار مفیدتری است تا برای «تولید» آن عدد.

\section{تحلیل خطا}
خطاهای مشاهده‌شده در پنج دسته جای می‌گیرند:
\begin{enumerate}
\item \textbf{دقت عددی:} منبع درست است، ولی مدل عدد گردشده‌ی مقدمه را به‌جای عدد دقیق جدول می‌آورد (پرسش \lr{h07}).
\item \textbf{ترجمه:} ابهام واژه‌های فارسی مانند «چند»، و ترجمه‌ی تحت‌اللفظی اصطلاحات (که با واژه‌نامه کاهش یافت).
\item \textbf{کیفیت متن منبع:} جدول‌های تکه‌شده (۴۱ درصد جدول‌ها هنگام بخش‌بندی دو یا چند تکه شده بودند)، اسکن‌های قدیمی و معادله‌های تصویری.
\item \textbf{ارزیابی ناقص:} در پرسش‌های \lr{ri}، مقاله‌های دیگری که همان داده را دارند غلط شمرده می‌شوند. بازبینی نشان داد که بخش درخور توجهی از «خطا»های رتبه‌ی اول در این گروه در واقع مقاله‌ی مرتبط‌اند.
\item \textbf{نگهبان «پیدا نشد»:} در یک ارزیابی پیشین، ۹ پرسش از ۴۰ پرسش بی‌جواب رد نشدند و سامانه برایشان منبع نامرتبط نشان داد.
\end{enumerate}
از این پنج دسته، دسته‌ی اول و دوم به مدل زبانی، دسته‌ی سوم به مرحله‌ی تبدیل سند، دسته‌ی چهارم به روش ارزیابی و دسته‌ی پنجم به آستانه‌های نگهبان مربوط‌اند. پس بهبود سامانه به یک جزء خاص محدود نیست و هر جزء سهم جداگانه‌ای در خطای نهایی دارد.

% ======================================== فصل ۶ ======================================================
\chapter{بحث، محدودیت‌ها و نتیجه‌گیری}

\section{پاسخ به پرسش‌های پژوهش}
\textbf{پرسش ۱ (\lr{RAG} در برابر مدل بدون منبع).} در مجموعه‌ی آزمایشی \lr{HITRAN}، \lr{RAG} دقت پاسخ عددی مدل \lr{Qwen3-4B} را از صفر به ۱۴ پرسش از ۱۵ پرسش رساند و امتناع درست در پرسش بی‌جواب حفظ شد. این نتیجه روی مجموعه‌ی کوچکی به دست آمده است، اما جهت اثر روشن است: برای عددهای کم‌تکرار ادبیات تخصصی، مدل کوچک بدون منبع کاربردی نیست. تنها خطای باقی‌مانده نیز از نوع خواندن بود و نه بازیابی، که نشان می‌دهد در این مجموعه گلوگاه اصلی، دقت مدل در انتخاب عدد درست از میان بخش‌های بازیابی‌شده است.

\textbf{پرسش ۲ (روش بازیابی).} هیچ روش واحدی در همه‌ی حوزه‌ها برتر نبود. در \lr{HITRAN} ترکیب برداری و \lr{BM25} بهترین بود و در مجموعه‌ی کوچک اپتیک \lr{BM25} به‌تنهایی. در مجموعه‌ی بزرگ، جست‌وجوی ترکیبی پایه‌ی سامانه ماند و تنها بهبود معنادار آماری از رتبه‌بندی دوم با مدل داور به دست آمد (۰/۵۱۸ به ۰/۵۹۰، با بهبود جفتی ۰/۰۲۷ تا ۰/۱۱۵ در بازه‌ی اطمینان ۹۵ درصد). رتبه‌بند متقاطع عمومی \lr{bge-reranker} نتیجه را بدتر کرد. این یافته با این دیدگاه سازگار است که مدل‌های آموزش‌دیده روی متن عمومی بدون آزمون مستقل به متن علمی تخصصی تعمیم نمی‌یابند.

\textbf{پرسش ۳ (پرسش فارسی).} مسیر ترجمه برای جست‌وجوی برداری بی‌هزینه بود، ولی \lr{BM25} و در نتیجه جست‌وجوی ترکیبی را تضعیف کرد (\lr{MRR} ترکیبی از ۰/۷۹ به ۰/۷۲). واژه‌نامه‌ی ثابت اصطلاحات فیزیک بخشی از خطاهای ترجمه را رفع کرد.

\textbf{پرسش ۴ (استخراج خودکار).} پاسخ منفی است. در ۲۴ مقاله‌ی آزموده، هیچ فرمول پاشندگی‌ای که مدل زبانی به‌طور خودکار استخراج کرد تا دقت ۰/۰۱ درست نبود. داده‌ی دقیق تنها با استخراج نیمه‌دستی و آزمون با عدد خود مقاله به دست آمد. سه بررسی خودکار (عدد در متن، شاهد عینی و بازه‌ی فیزیکی) جلوی عددسازی آشکار را می‌گیرند، ولی خطای خواندن ضرایب را نمی‌گیرند، زیرا ضریبی که از ستون اشتباه خوانده شده هم «در متن منبع هست» و هم در بازه‌ی معقول قرار دارد.

\section{محدودیت‌ها}\label{sec:limits}
\subsection*{ارزیابی}
\begin{itemize}
\item بازه‌ی اطمینان \lr{MRR} در ارزیابی ۲۰۰ پرسشی حدود $\pm0.06$ است. تفاوت‌های کوچک‌تر از این (مانند اثر پیش‌واکشی بلور) قابل‌اتکا نیستند.
\item تنظیم سامانه و گزارش نتیجه روی یک مجموعه‌ی پرسش انجام شد و مجموعه‌ی آزمون جداگانه‌ای کنار گذاشته نشد. پس عددهای گزارش‌شده احتمالاً کمی خوش‌بینانه‌اند.
\item داوری مرتبط بودن کامل نیست: هر پرسش فقط یک یا چند مقاله‌ی «طلایی» دارد، و مقاله‌های دیگری که همان داده را دارند غلط شمرده می‌شوند. برای داوری نهایی، داوری درجه‌بندی‌شده (همه‌ی مقاله‌هایی که به پرسش پاسخ می‌دهند) لازم است.
\item گروه‌های ترمواپتیک، ضرایب $d$ و مقاله‌ی فرمول فقط ۳ تا ۷ پرسش دارند.
\item پاسخ نوشتاری فقط روی مجموعه‌ی آزمایشی \lr{HITRAN} (۱۶ پرسش) ارزیابی شد. ارزیابی وفاداری پاسخ با \lr{RAGAS}، مقایسه‌ی دو مدل برداری و اثر اندازه‌ی بخش‌ها، که در برنامه‌ی اولیه بود، انجام نشد.
\item ارزیابی با کاربر واقعی انجام نشده و سناریوهای مشاور بلور ساخته‌ی خود نگارنده‌اند.
\end{itemize}

\subsection*{داده و مجموعه}
\begin{itemize}
\item ۵۹۸ سند (۲۱ درصد) فقط چکیده دارند و ۳۳ سند فقط مشخصات. حدود ۴۳ فایل \lr{PDF} مقاله‌ی دیگری را در خود دارند و ۴۷ مقاله‌ی اسکن‌شده‌ی قدیمی متن خراب دارند.
\item مدل ترمواپتیک فقط برای ۶۷ بلور از ۱۸۶ بلور محاسبه‌گر وجود دارد و تانسور $d$ برای چند بلور ناقص است.
\item مجموعه به‌روزرسانی خودکار ندارد و فقط با دخالت دستی بزرگ می‌شود.
\end{itemize}

\subsection*{روش و زیرساخت}
\begin{itemize}
\item مدل برداری فقط انگلیسی است و پرسش فارسی از مسیر ترجمه می‌گذرد.
\item محاسبه‌ی تطبیق فاز (به‌جز بخش \lr{SPDC}) در تقریب موج تخت است. بازده تبدیل با کانونی‌سازی، انتشار تپ و اثرهای گرمایی مدل نشده‌اند.
\item رتبه‌بندی دوم و بخشی از برچسب‌زنی به اعتبار رایگان سرویس‌های بیرونی وابسته است. بدون آن، کیفیت جست‌وجو به حالت ترکیبی برمی‌گردد. پردازش سنگین هم در \lr{Colab} رایگان انجام شد که پایدار نیست.
\item به دلیل حق نشر متن مقاله‌ها، سایت عمومی نیست و تکرار کامل نتایج به دسترسی به همان مجموعه‌ی \lr{PDF} نیاز دارد. کد، پرسش‌های مرجع، گزارش‌های ارزیابی و داده‌ی فیزیکی استخراج‌شده در مخزن عمومی هست.
\end{itemize}

\section{کارهای آینده}
\begin{enumerate}
\item ساخت مجموعه‌ی ارزیابی بزرگ‌تر با داوری درجه‌بندی‌شده، دست‌کم ۳۰ پرسش برای هر گروه داده، و یک مجموعه‌ی آزمون جدا که در تنظیم سامانه استفاده نشود.
\item استخراج نیمه‌خودکار با تأیید انسانی: مدل زبانی جدول‌های حاوی \lr{pm/V} و $dn/dT$ را پیشنهاد کند و انسان تأیید کند، با همان آزمون «عدد مقاله» به‌عنوان دروازه‌ی ورود.
\item استفاده از یک مدل برداری چندزبانه برای پرسش فارسی بدون واسطه‌ی ترجمه.
\item گسترش مدل فیزیکی: بازده تبدیل با کانونی‌سازی \lr{Boyd-Kleinman} \cite{boydkleinman} و پهنای پذیرش تپی.
\item آزمودن مشاور بلور با سناریوهای واقعی آزمایشگاه‌های دانشگاه.
\item تکمیل مسیر \lr{HITRAN}: ساخت مجموعه‌ی مرجع نهایی با ۵۰ تا ۱۰۰ نمونه‌ی بررسی‌شده و یکسان‌سازی واحدها.
\end{enumerate}

\section{نتیجه‌گیری}
در این پروژه یک سامانه‌ی \lr{RAG} برای بازیابی اطلاعات و استخراج خواص فیزیکی مواد از مقاله‌های علمی طراحی، پیاده‌سازی و ارزیابی شد. سامانه از مقاله‌ی \lr{PDF} تا پاسخ مستند با شماره‌ی صفحه و محل روی صفحه کار می‌کند، پرسش فارسی را می‌پذیرد و در نبود شاهد صریحاً «پیدا نشد» می‌گوید. روی مجموعه‌ای با ۲۸۴۵ سند و ۶۸۱۹۰ بخش، جست‌وجوی ترکیبی با رتبه‌بندی دوم در ۷۴/۵ درصد پرسش‌ها مقاله‌ی طلایی را در ۱۰ نتیجه‌ی اول آورد.

مهم‌ترین درس پروژه این است که در کار فیزیکی، \lr{RAG} ابزار \textbf{یافتن منبع} است و نه جایگزین خواندن آن. مدل زبانی منبع درست را خوب پیدا می‌کند، ولی ضرایب یک فرمول را به‌طور قابل‌اتکا نمی‌خواند. ترکیب بازیابی خودکار با آزمون هر عدد در برابر مقدار گزارش‌شده در خود مقاله، پایگاه داده‌ای ساخت که محاسبه‌هایش با ادبیات هم‌خوان است و در مسیر ساختش چند ناسازگاری در خود ادبیات هم آشکار شد. این الگو، یعنی «بازیابی خودکار، خواندن انسانی، آزمون خودکار»، به حوزه‌های دیگر فیزیک که داده‌ی پایه‌شان در مقاله‌ها پراکنده است نیز قابل‌انتقال است.

% ======================================== پیوست‌ها ===================================================
\appendix
\chapter{واژه‌نامه}
\begin{table}[htbp]
\centering
\caption{واژه‌نامه‌ی فارسی به انگلیسی اصطلاحات گزارش.}
\label{tab:glossary}
\renewcommand{\arraystretch}{1.25}
\small
\begin{tabular}{p{4.2cm}p{5.4cm}p{5.0cm}}
\toprule
فارسی & انگلیسی & توضیح \\
\midrule
تولید تقویت‌شده با بازیابی & \lr{Retrieval-Augmented Generation (RAG)} & پاسخ مدل زبانی بر پایه‌ی متن منابع بازیابی‌شده \\
پایگاه داده‌ی برداری & \lr{Vector database} & ذخیره و جست‌وجوی سریع بردارهای معنایی \\
نمایش برداری، بردارسازی & \lr{Embedding} & نگاشت متن به بردار عددی \\
بخش & \lr{Chunk} & واحد جست‌وجو: یک پاراگراف، جدول یا تکه‌ی جدول \\
بازیابی ترکیبی & \lr{Hybrid retrieval} & ترکیب جست‌وجوی برداری و کلیدواژه‌ای \\
ادغام معکوس رتبه & \lr{Reciprocal Rank Fusion (RRF)} & رابطه‌ی (\ref{eq:rrf}) \\
رتبه‌بندی دوم & \lr{Re-ranking} & مرتب‌سازی دوباره‌ی نامزدهای اول با مدل دقیق‌تر \\
توهم & \lr{Hallucination} & تولید محتوای باورپذیر ولی نادرست توسط مدل زبانی \\
نویسه‌خوانی نوری & \lr{OCR} & تبدیل تصویر صفحه به متن \\
فرمول پاشندگی (سلمایر) & \lr{Dispersion (Sellmeier) formula} & رابطه‌ی (\ref{eq:sellmeier}) \\
تطبیق فاز & \lr{Phase matching} & رابطه‌ی (\ref{eq:dk}) \\
تطبیق فاز شبه‌فازی & \lr{Quasi-phase matching (QPM)} & رابطه‌ی (\ref{eq:qpm}) \\
جدایی مسیر & \lr{Walk-off} & رابطه‌ی (\ref{eq:walkoff}) \\
فروتبدیل پارامتری خودبه‌خودی & \lr{SPDC} & منبع جفت فوتون درهم‌تنیده \\
میانگین معکوس رتبه & \lr{Mean Reciprocal Rank (MRR)} & رابطه‌ی (\ref{eq:metrics}) \\
\bottomrule
\end{tabular}
\end{table}
اصطلاحات اصلی گزارش در جدول \ref{tab:glossary} آمده‌اند. همین واژه‌نامه، با افزودن اصطلاحات بیشتر، در ترجمه‌ی پرسش‌های فارسی به مدل زبانی داده می‌شود (بخش ۳-۵-۳).

\chapter{منشأ عددهای گزارش}
هر عدد این گزارش از یکی از فایل‌های زیر در مخزن کد پروژه آمده است:
\begin{itemize}
\item ترکیب مجموعه، بخش‌ها و چارک مجله‌ها: شمارش مستقیم پایگاه \lr{Qdrant} و تطبیق مشخصات \lr{OpenAlex} با فهرست \lr{SCImago}، ۳ آبان ۱۴۰۵ (\lr{\texttt{docs/\allowbreak matrag\_report\_aban1405.html}}).
\item آزمون‌های اولیه: \lr{\texttt{results/\allowbreak eval/\allowbreak }} در مخزن کد.
\item ارزیابی ۲۰۰ پرسشی: \lr{\texttt{scripts/\allowbreak quality/\allowbreak reports/\allowbreak search-2026-10-26.md}} و فایل \lr{\texttt{.json}} همراه آن. بازه‌های اطمینان با \lr{bootstrap} روی همین فایل محاسبه شد.
\item رتبه‌بند \lr{bge}: \lr{\texttt{docs/\allowbreak reranker\_eval.md}}.
\item داده‌ی فیزیکی: \lr{\texttt{data/\allowbreak rag-optics/\allowbreak nonlinear\_coefficients.yml}}، \lr{\texttt{thermo\_optic.yml}}، \lr{\texttt{paper\_formulas.yml}} و \lr{\texttt{crystal\_properties.yml}}.
\item مثال‌های حل‌شده‌ی بخش ۵-۵ و شکل‌های \ref{fig:bbo-n}، \ref{fig:bbo-acc} و \ref{fig:ppln}: برنامه‌ی \lr{\texttt{scripts/\allowbreak calc\_figures.py}} همراه این پایان‌نامه. شکل \ref{fig:mrr} با همین برنامه از عددهای جدول \ref{tab:groups} رسم شده است.
\end{itemize}

\chapter{اجرای سامانه}
فرمان‌های زیر سامانه را نصب، ارزیابی و اجرا می‌کنند. آزمون‌ها بدون اینترنت و بدون مدل زبانی اجرا می‌شوند.
\codebox{\texttt{%
pip install -e ".[dev]"\\
matrag -{}-corpus hitran add 1906.01475\hfill \# add a paper from arXiv\\
matrag -{}-corpus hitran evaluate retrieval\hfill \# Recall@k and MRR, three retrieval modes\\
matrag -{}-corpus hitran evaluate answers\hfill \# RAG vs. LLM alone, "not found" test\\
matrag reference hitran results.csv -{}-molecule CO2 -{}-nu-min 5007.7 -{}-nu-max 5007.9\\
matrag app\hfill \# web UI at http://127.0.0.1:7860\\
pytest\hfill \# tests (offline, no LLM)}}

عددهای مثال‌های حل‌شده و نمودارهای فصل ۵ با فرمان زیر دوباره ساخته می‌شوند:
\codebox{\texttt{python3 scripts/\allowbreak calc\_figures.py}}

% ======================================== منابع ======================================================
\chapter*{منابع و مآخذ}
\addcontentsline{toc}{chapter}{منابع و مآخذ}
\markboth{منابع و مآخذ}{منابع و مآخذ}
\begin{latin}
\renewcommand{\bibname}{\rl{منابع و مآخذ}}
\renewcommand{\chapter}[2]{}
\begin{thebibliography}{99}
\bibitem{ji} Z. Ji \textit{et al.}, ``Survey of hallucination in natural language generation,'' \textit{ACM Comput. Surv.}, vol. 55, no. 12, art. 248, 2023.
\bibitem{lewis} P. Lewis \textit{et al.}, ``Retrieval-augmented generation for knowledge-intensive NLP tasks,'' in \textit{Advances in Neural Information Processing Systems}, vol. 33, pp. 9459--9474, 2020.
\bibitem{polyanskiy} M. N. Polyanskiy, ``Refractiveindex.info database of optical constants,'' \textit{Sci. Data}, vol. 11, art. 94, 2024.
\bibitem{hitran} I. E. Gordon \textit{et al.}, ``The HITRAN2020 molecular spectroscopic database,'' \textit{J. Quant. Spectrosc. Radiat. Transf.}, vol. 277, art. 107949, 2022.
\bibitem{vaswani} A. Vaswani \textit{et al.}, ``Attention is all you need,'' in \textit{Advances in Neural Information Processing Systems}, vol. 30, pp. 5998--6008, 2017.
\bibitem{gao} Y. Gao \textit{et al.}, ``Retrieval-augmented generation for large language models: A survey,'' arXiv:2312.10997, 2023.
\bibitem{sbert} N. Reimers and I. Gurevych, ``Sentence-BERT: Sentence embeddings using Siamese BERT-networks,'' in \textit{Proc. EMNLP-IJCNLP}, pp. 3982--3992, 2019.
\bibitem{dpr} V. Karpukhin \textit{et al.}, ``Dense passage retrieval for open-domain question answering,'' in \textit{Proc. EMNLP}, pp. 6769--6781, 2020.
\bibitem{bge} S. Xiao, Z. Liu, P. Zhang, and N. Muennighoff, ``C-Pack: Packaged resources to advance general Chinese embedding,'' arXiv:2309.07597, 2023.
\bibitem{hnsw} Y. A. Malkov and D. A. Yashunin, ``Efficient and robust approximate nearest neighbor search using hierarchical navigable small world graphs,'' \textit{IEEE Trans. Pattern Anal. Mach. Intell.}, vol. 42, no. 4, pp. 824--836, 2020.
\bibitem{bm25} S. Robertson and H. Zaragoza, ``The probabilistic relevance framework: BM25 and beyond,'' \textit{Found. Trends Inf. Retr.}, vol. 3, no. 4, pp. 333--389, 2009.
\bibitem{rrf} G. V. Cormack, C. L. A. Clarke, and S. B\"uttcher, ``Reciprocal rank fusion outperforms Condorcet and individual rank learning methods,'' in \textit{Proc. 32nd ACM SIGIR}, pp. 758--759, 2009.
\bibitem{nogueira} R. Nogueira and K. Cho, ``Passage re-ranking with BERT,'' arXiv:1901.04085, 2019.
\bibitem{yariv} A. Yariv and P. Yeh, \textit{Optical Waves in Crystals}. New York: Wiley, 1984.
\bibitem{boyd} R. W. Boyd, \textit{Nonlinear Optics}, 3rd ed. Burlington, MA: Academic Press, 2008.
\bibitem{armstrong} J. A. Armstrong, N. Bloembergen, J. Ducuing, and P. S. Pershan, ``Interactions between light waves in a nonlinear dielectric,'' \textit{Phys. Rev.}, vol. 127, no. 6, pp. 1918--1939, 1962.
\bibitem{boydkleinman} G. D. Boyd and D. A. Kleinman, ``Parametric interaction of focused Gaussian light beams,'' \textit{J. Appl. Phys.}, vol. 39, no. 8, pp. 3597--3639, 1968.
\bibitem{bennink} R. S. Bennink, ``Optimal collinear Gaussian beams for spontaneous parametric down-conversion,'' \textit{Phys. Rev. A}, vol. 81, art. 053805, 2010.
\bibitem{cde} M. C. Swain and J. M. Cole, ``ChemDataExtractor: A toolkit for automated extraction of chemical information from the scientific literature,'' \textit{J. Chem. Inf. Model.}, vol. 56, no. 10, pp. 1894--1904, 2016.
\bibitem{mp} A. Jain \textit{et al.}, ``Commentary: The Materials Project: A materials genome approach to accelerating materials innovation,'' \textit{APL Mater.}, vol. 1, art. 011002, 2013.
\bibitem{dmitriev} V. G. Dmitriev, G. G. Gurzadyan, and D. N. Nikogosyan, \textit{Handbook of Nonlinear Optical Crystals}, 3rd ed. Berlin: Springer, 1999.
\bibitem{nikogosyan} D. N. Nikogosyan, \textit{Nonlinear Optical Crystals: A Complete Survey}. New York: Springer, 2005.
\bibitem{paperqa} M. D. Skarlinski \textit{et al.}, ``Language agents achieve superhuman synthesis of scientific knowledge,'' arXiv:2409.13740, 2024.
\bibitem{ragas} S. Es, J. James, L. Espinosa-Anke, and S. Schockaert, ``RAGAS: Automated evaluation of retrieval augmented generation,'' arXiv:2309.15217, 2023.
\bibitem{hapi} R. V. Kochanov \textit{et al.}, ``HITRAN Application Programming Interface (HAPI): A comprehensive approach to working with spectroscopic data,'' \textit{J. Quant. Spectrosc. Radiat. Transf.}, vol. 177, pp. 15--30, 2016.
\bibitem{docling} C. Auer \textit{et al.}, ``Docling technical report,'' arXiv:2408.09869, 2024.
\bibitem{qwen3} A. Yang \textit{et al.} (Qwen Team), ``Qwen3 technical report,'' arXiv:2505.09388, 2025.
\bibitem{roberts} D. A. Roberts, ``Simplified characterization of uniaxial and biaxial nonlinear optical crystals: A plea for standardization of nomenclature and conventions,'' \textit{IEEE J. Quantum Electron.}, vol. 28, no. 10, pp. 2057--2074, 1992.
\bibitem{efron} B. Efron and R. J. Tibshirani, \textit{An Introduction to the Bootstrap}. New York: Chapman \& Hall, 1993.
\bibitem{fedrizzi} A. Fedrizzi, T. Herbst, A. Poppe, T. Jennewein, and A. Zeilinger, ``A wavelength-tunable fiber-coupled source of narrowband entangled photons,'' \textit{Opt. Express}, vol. 15, no. 23, pp. 15377--15386, 2007.
\bibitem{eimerl} D. Eimerl, L. Davis, S. Velsko, E. K. Graham, and A. Zalkin, ``Optical, mechanical, and thermal properties of barium borate,'' \textit{J. Appl. Phys.}, vol. 62, no. 5, pp. 1968--1983, 1987.
\bibitem{gayer} O. Gayer, Z. Sacks, E. Galun, and A. Arie, ``Temperature and wavelength dependent refractive index equations for MgO-doped congruent and stoichiometric LiNbO$_3$,'' \textit{Appl. Phys. B}, vol. 91, pp. 343--348, 2008.
\end{thebibliography}
\end{latin}

% ======================================== چکیده‌ی انگلیسی ==========================================
\begin{latin}
\clearpage
\thispagestyle{plain}
\begin{center}
{\Large\bfseries Abstract}
\end{center}
\addcontentsline{toc}{chapter}{\rl{چکیده‌ی انگلیسی}}

Basic materials data such as refractive-index dispersion formulas, second-order nonlinear coefficients, thermo-optic coefficients and spectral line parameters are scattered across the tables and equations of thousands of papers, and general-purpose large language models often invent such numbers without a verifiable source. This project designed and implemented a retrieval-augmented generation (RAG) system that converts PDF papers into structured text and tables, indexes them in vector databases (ChromaDB locally, Qdrant for the web site), and answers Persian or English questions with a citation to the paper, page and location on the page. Retrieval is hybrid: dense search with bge-base-en-v1.5 embeddings and BM25 keyword search fused by reciprocal rank fusion. Structured property extraction comes with three automatic checks: the value is present in the source, the quoted evidence is verbatim, and the value lies in a physically plausible range.

The system was first tested on molecular spectroscopy with HITRAN as ground truth and then extended to linear and nonlinear optical crystals, with a final collection of 2,845 documents, 2,214 of them in full text, split into 68,190 searchable passages. On the HITRAN pilot set, RAG with a local Qwen3-4B model gave the correct value for 14 of 15 questions, while the same model without sources gave none. On 200 gold questions over the large collection, MRR@10 was 0.518 for hybrid search and 0.590 after LLM re-ranking of the top 30 (paired 95\% bootstrap interval of the gain: +0.027 to +0.115). A tested physical database and calculators for phase matching and photon-pair generation were built on top; their outputs agree with values reported in the original papers, for example a type-I SHG phase-matching angle of 22.83$^\circ$ in BBO at 1.064~$\mu$m against 22.8$^\circ$ in the literature. A key negative result is that fully automatic extraction of formula coefficients by a language model did not reach the needed accuracy; accurate data required semi-manual extraction checked against each paper's own reported values.

\vspace{6pt}
\noindent\textbf{Keywords:} retrieval-augmented generation, vector database, hybrid retrieval, materials property extraction, refractive index, nonlinear optics, phase matching

\clearpage
\thispagestyle{plain}
\vspace*{4cm}
\begin{center}
\textbf{University of Isfahan}\\
Faculty of Physics\\
B.Sc. Physics Project Report (3 credits)\\[10pt]
{\large\bfseries Design and Implementation of an Intelligent Information Retrieval and Materials Property Extraction System for Scientific Papers Using RAG Architecture and Vector Databases}\\[10pt]
Student: Amin Hossein Sadidi (Student ID: 4024213027)\\
Supervisor: Dr. Mohammadi\\
November 2026
\end{center}
\end{latin}

\end{document}
